\documentclass[12pt,a4paper,onecolumn]{article}
\title{Digest of The New York Times}
\author{The New York Times}
% \definecolor{dark-green}{rgb}{0,0.25,0}
% \usepackage[utf8]{inputenc}

\usepackage{fontspec,xunicode,xltxtra}
\setmainfont[Ligatures=Required,Mapping=tex-text]{Georgia}
% \setmainfont[Ligatures=Required,Mapping=tex-text]{Palatino Linotype}
\setsansfont[Ligatures=Required,Mapping=tex-text]{Myriad Pro}
\setmonofont[Ligatures=Required,Mapping=tex-text]{Courier New}
\usepackage{xeCJK}
\setCJKmainfont{宋体}
\setCJKmonofont{黑体}

% \usepackage{indentfirst}


\usepackage[dvips,dvipsnames,svgnames]{xcolor}
\definecolor{light-gray}{gray}{0.95}

%生成PDF的链接
\usepackage[dvipdfmx]{hyperref}
\hypersetup{
    bookmarks=true,         % show bookmarks bar?
    bookmarksopen=true,
    pdfpagemode=UseNone,    % options: UseNode, UseThumbs, UseOutlines, FullScreen
    pdfstartview=FitH,
    pdfborder=1,
    pdfhighlight=/P,
    pdfauthor={wuxch},
    unicode=true,           % non-Latin characters in Acrobat’s bookmarks
    colorlinks,             % false: boxed links; true: colored links
    linkcolor=blue,         % color of internal links
    citecolor=blue,        % color of links to bibliography
    filecolor=magenta,      % color of file links
    urlcolor=cyan           % color of external links
}
% 使用目录
\usepackage{makeidx}
% 使用图形
\usepackage{graphicx}


\usepackage{titletoc}
\dottedcontents{section}[1.5em]{}{2.3em}{1pc}

% % 页眉页脚的设定
% \usepackage{fancyhdr}
% \pagestyle{fancy}
% \lhead{\includegraphics[height=0.2in]{zte_logo}}

% \usepackage{footmisc}
\usepackage{pfnote}
\usepackage{fnpos}

% % 嵌入的代码显示
% \usepackage{listings}
% \lstset{language=C++, breaklines, extendedchars=false}
% \lstset{keywordstyle=\color{blue},
%         commentstyle=\color{SeaGreen},
%         stringstyle=\ttfamily,
%         showstringspaces=false,
%         backgroundcolor=\color{light-gray}}
% \lstset{basicstyle=\small}

\usepackage{booktabs}

\usepackage{tipa}
\usepackage{phonetic}
\usepackage[gen]{eurosym}

% \usepackage{CJK}
\usepackage{makeidx}
% \usepackage{glossary}

\usepackage{fancyhdr}

\usepackage{lastpage}
\cfoot{\textsf{\thepage\ of \pageref*{LastPage}}}

\pagestyle{fancy}
\fancyhead{}
\renewcommand{\headrulewidth}{0pt}
\pagenumbering{arabic}

\usepackage{graphicx}
\usepackage{picinpar}

\setlength{\parindent}{0in}
\setlength{\parskip}{2.5ex plus4mm minus3mm}
% \setlength{\parskip}{.3\baselineskip}

\setcounter{section}{0}

% bibtex
\usepackage[super]{natbib}
% \usepackage[globalcitecopy,labelstoglobalaux]{bibunits}
% \makeatletter
% \renewcommand{\@bibunitname}{\jobname.\the\@bibunitauxcnt}
% \makeatother

% \renewcommand{\refname}{缩略语}
% \renewcommand{\refname}{}

\makeatletter                   % Reference list option change
\renewcommand\@biblabel[1]{#1:  } % from [1] to 1: 
\makeatother                    %

\newcommand{\citenumfont}[1]{\textsf{#1}}

% \usepackage{tipa}
% \usepackage{phonetic}

\newcommand{\phonetic}[1]{{\fontspec[Colour=CC00CC]{Microsoft Sans Serif} #1}}


%%% Local Variables:
%%% TeX-master: "medical_trivia"
%%% End:

\makeindex
\begin{document}
\begin{CJK*}{GBK}{song}
\date{}
% \maketitle
\tableofcontents
\pagebreak
\section{Treasury Secretary Facing a Defining Moment}

WASHINGTON -- All three of President Obama's top economic advisers were on message when they
appeared Sunday on separate television talk shows. Treasury Secretary Timothy F. Geithner, they
said, had concluded, based on lawyers' advice, that he could not stop the \$165 million in bonuses
that the American International Group was even then doling out to hundreds of employees.

But when Mr.~Geithner and other officials met at the White House that night, the president's
political advisers -- who had agreed to the day's message -- decided the growing outcry left
Mr.~Obama no choice but to publicly second-guess his Treasury secretary.

The next morning on camera, the president said he had directed Mr.~Geithner to find a legal way ``to
block these bonuses and make the American taxpayers whole.''

Thus began perhaps the worst week in a string of bad weeks for the Treasury secretary. The mixed
messages on A.I.G. gave further ammunition to critics who had begun questioning Mr.~Geithner's
credibility as the administration's point man on the economy, an essential commodity if he is to
help restore consumer confidence. Fair or not, questions about why Mr.~Geithner did not know sooner
about the A.I.G. bonuses and act to stop them threaten to overwhelm his achievements and undermine
Mr.~Obama's overall economic agenda. Edward M. Liddy, chief executive of A.I.G., told Congress on
Wednesday that he generally deals with Fed officials, figuring they would keep Treasury informed.
Even if they had, administration officials say that Mr.~Geithner's instincts are that government
should not dictate compensation issues to businesses.

The controversy comes as Mr.~Geithner is about to announce details of the restructured bank rescue
program, and it clouds prospects for more rescue funds that the administration is all but certain to
need.

The once-heralded Wall Street credentials of Mr.~Geithner, formerly the president of the Federal
Reserve Bank of New York, were already marred by false starts in revamping the Bush administration's
bank rescue program, even as his perceived closeness to financiers and unease with populist politics
left Main Street skeptical.

On Wednesday, a junior Republican in Congress and some traders on Wall Street went so far as to call
for him to quit or be fired. The Republican leader of the House, Representative John A. Boehner of
Ohio, told a conservative talk-radio host that the secretary is ``on thin ice.''

But Mr.~Geithner's boss, the president, interjected a vote of ``complete confidence.''

``Tim Geithner didn't draft these contracts with A.I.G.,'' Mr.~Obama told reporters on the White
House lawn as he left for California on Wednesday. ``There has never been a secretary of the
Treasury, except maybe Alexander Hamilton right after the Revolutionary War, who's had to deal with
the multiplicity of issues that Secretary Geithner is having to deal with -- all at the same time.''

``He is making all the right moves in terms of playing a bad hand,'' the president continued. ``And
what we need to be doing is making sure that we are providing him the support that he needs.''

Mr.~Geithner is shouldering more crises on his slight frame than most Treasury secretaries ever
have. And he is doing so without the usual complement of Treasury assistants because of
administration delays in vetting potential nominees -- a consequence in part of its efforts to avoid
embarrassments like the disclosures of Mr.~Geithner's past tax lapses, which nearly doomed his
nomination.

Since before his confirmation in late January, Mr.~Geithner has juggled a crushing workload:
overhauling the Bush administration's discredited financial bailout program; helping with
Mr.~Obama's nearly \$800 billion economic stimulus plan; and managing the government effort to
salvage the auto industry.

Mr.~Geithner is now fashioning a new federal regulatory structure for the financial industry to
replace the one that failed. He has developed a housing program that aims to avert up to nine
million more foreclosures, and programs for getting credit flowing to small businesses and consumers
as well as the major financial giants.

At 47, the same age as the president, Mr.~Geithner works out at 5:30 a.m., gets to his desk by 6:30
and leaves 15 hours later.

On Tuesday last week, as he prepared for a meeting in London of the finance ministers of the Group
of 20 nations, Mr.~Geithner learned that A.I.G. by Sunday would send out the bonuses to employees at
its financial products unit, which developed the risky derivatives now blamed for the global credit
crisis.

With few senior political appointees on hand, the word came from one of the numerous career civil
servants who keep the Treasury functioning through changes of administration, according to an
official.

Mr.~Geithner consulted lawyers. They told him the government could not override the contracts that
the insurance conglomerate had signed in early 2008, when its financial products unit already was
fast losing money.

On Wednesday evening, Mr.~Geithner called Mr.~Liddy, and demanded that he renegotiate payments. The
next morning, Mr.~Geithner informed White House advisers. Later that day a senior adviser, David
Axelrod, informed the president.

On Friday, Mr.~Liddy said he could not block the bonuses; he did agree to reduce future executive
bonuses set for July 15 and Sept.~15. With Mr.~Geithner in London, Treasury officials tried to
manage the potential criticism by leaking word to selected news media on Saturday. On Sunday, the
economic advisers went on TV.

The A.I.G. tempest has been especially explosive for Mr.~Geithner because, as president of the New
York Fed, he was the one administration official who had been involved in the Bush-era bailouts.

Once A.I.G. was under the Fed's control, its compensation plans hardly came up, according to
officials. In December, an initial \$55 million in bonuses went out with hardly a stir.

\section{Chinese City Bolsters Scant Consumer Spending With Free Vouchers}

HANGZHOU, China -- Shou Xuejun was a popular man one afternoon early this month. As a finance
official in this town of ancient canals and textile factories, he got to hand out shopping vouchers
worth nearly \$30 apiece to 20 elderly residents.

Chen Jue, 87, could hardly believe her good luck. ``Can these vouchers really buy things?'' she
asked, amazed, at a community center decorated with a six-foot-high photograph of Mao. ``The
government is very good to us! I will buy flowers, and I will also buy goldfish.'' Since January,
more than a fifth of the five million residents of the Hangzhou area have received the same
vouchers. Hundreds of thousands more vouchers will be issued soon. Hangzhou officials hope the
coupons will stimulate buying and avert the worst of an economic downturn that has eviscerated
factory after factory here.

But even Mr.~Shou, who helps oversee the \$100 million program, acknowledges that vouchers will not
cure Hangzhou's slump -- or China's. The problem is more basic: China's economy thrives on making
things that overseas customers want to buy. Chinese consumers tend to be less eager shoppers.
Conditioned to save for medical expenses and old age, they sock away more than one-fourth of their
incomes -- far more than Westerners, especially Americans.

``People do not dare to spend more because of the underdevelopment of the social welfare system,''
Mr.~Shou said. ``That problem cannot be solved in the short term.''

China's leaders have talked for years about the need to create a social safety net that would ease
the financial pressure on average Chinese and make consumer demand and the service sector a bigger
source of growth. Now that exports are slowing in the face of a global economic crisis, many
economists say, it is past time for action.

Despite recent increases in social spending, especially on education, the government has mostly
concentrated on building up industry and infrastructure, subsidizing a host of manufacturing costs.
That economic model has delivered world-beating growth for three decades and lifted hundreds of
millions out of poverty.

Chinese consumers have played a comparatively small part. Their spending accounts for just 35
percent of China's gross domestic product, compared with more than two-thirds for consumers in the
United States.

They save for good reason: at least one-fourth of the population has no health insurance at all,
according to official estimates. Hundreds of millions of others face crippling bills for treatment
of serious illnesses that are not covered by rudimentary insurance programs. Public pensions cover
less than one-third of workers. An estimated 130 million migrant workers are not protected by
unemployment insurance. Payments to the poor reach only a fraction of those eligible, according to
the China Development Research Foundation, a nonprofit group.

Prime Minister Wen Jiabao said this month that China must make domestic demand ``a basic point of
departure for stimulating economic growth.'' But weaving a social safety net that gives average
Chinese the confidence to loosen their purse strings would take years. Meanwhile, the overseas
consumer is in full flight: In February, exports fell by a record 25.7 percent. Some economists warn
that exports might not revive until next year.

With the Communist Party's political legitimacy closely tied to economic growth, China's leaders are
trying to save and create jobs with a \$584 billion stimulus plan tilted heavily toward
infrastructure. Outside economists say they hope the money goes to build schools or low-cost housing
for migrant families, rather than more highways and seaports to benefit industry.

The government has also pledged more social spending, led by a \$123 billion three-year initiative
to deliver basic, universal health care and health insurance to 9 in 10 Chinese. That follows a
three-year drive to provide compulsory, free education to students through ninth grade.

``This is step one of a 100-step process, but it is a very important step,'' Arthur Kroeber,
managing director of Dragonomics, a Beijing-based economic research firm, said of the health care
plan. ``There is political will, and they are putting real money on the table at a time when it is
not easy to do so.''

But in a January meeting with Finance Ministry officials, outside economists argued that the
government should do more now to protect the vulnerable. ``The government could afford to expand all
these programs quickly,'' David Dollar, the World Bank's country director, wrote in a report on the
bank's Web site describing the meeting.

``They really need a sweeping program to set up social security,'' said Fred Hu, a Goldman Sachs
economist who has advised the government on policy. ``That hasn't happened. It's still piecemeal
measures.''

The coastal city of Hangzhou, just south of Shanghai, has been hit less hard than some cities in
China's manufacturing belt. Still, officials say industrial output has shrunk since late last year,
and migrant-worker centers are packed with job seekers. At one center last week, the arrival of a
foreigner -- mistaken for a possible employer -- caused such commotion that security guards came
running.

Wang Pengfei, 21, was one of hundreds in Hangzhou hoping to find work. Until December, he helped
fill orders for the Hangzhou Donghua Chain Group Company, which makes chains for conveyor belts,
tractors, motorcycles and other machinery. Michael Xiao, that company's head of exports, said that
Donghua normally exported four-fifths of its production. Now those orders have fallen by 50 or 60
percent. In December, the factory laid off 500 of its 3,000 workers.

Hangzhou officials are hoping to soften such blows by giving vouchers to everyone from middle school
students to tourists. Some vouchers are aimed at the needy, including laid-off workers; others are
more broadly directed toward consumers in the hope of spurring spending.

A red and yellow banner stretched above the doors to Century Mart, a supermarket, announces:
``'Welcome, shoppers. Come to our store and spend your voucher.''

``I would buy this anyway, but this relieves our burden,'' said Zhang Beili, a 62-year-old retired
government worker, as she cashed a \$3 voucher to buy snacks and soft drinks.

Although some other Chinese cities are also experimenting with vouchers, government officials said
there was no plan now to expand the program nationwide. Many economists question whether vouchers
can revive consumption, saying many recipients may use them instead of spending cash, not in
addition.

Still, in Taiwan, where the economy is more consumer-based, President Ma Ying-jeou said that a
program that gave every Taiwanese citizen vouchers worth the equivalent of \$102 was a success. In
an interview last month, Mr.~Ma said that 50,000 retailing jobs were saved and that about two-thirds
of 1 percent was added to the economic output.

Whether Hangzhou will be so lucky is unclear. At the least, Mr.~Shou said, the city is sending a
message to its citizens to ``have confidence'' and spend.

\section{36 Hours in Shanghai}

NOW that the Beijing Olympics are but a memory, the spotlight in China is moving to Shanghai as that
city gears up to host the 2010 World Expo. With an anticipated 70 million visitors and 200
participating countries, the six-month World's Fair will be enormous by any measure -- not that
Shanghai has ever needed an excuse to party. While the global economic slowdown has had its impact,
Beijing's naughty sister is still up to her tricks: from the flashing neon signs and light-bedazzled
skyscrapers to the throbbing clubs and houses from the foreign-concession era hiding their decadent
secrets. But beyond the clich\'es, mainland China's most cosmopolitan city still offers a breadth of
experiences.

\textbf{Friday}
\begin{description}
\item{7 p.m.} 1) JOURNEY TO KITSCH

Tonight is about embracing the kitsch. So set the tone by taking the Bund Sightseeing Tunnel, a
Disneyesque ride from the historic Bund area (look for the sign across from the Peace Hotel on
Nanjing Road East) to the futuristic Pudong district. Buy the 40 yuan ticket (about \$5.70 at 7 yuan
to the dollar), and a silver pod will shuttle you across the Huangpu River through an extravaganza
of pulsing, flashing and spiraling lights, creepy blow-up dolls and even creepier voice-overs
(``hell and paradise,'' ``nascent magma''). Don't ask questions; just sit back and look forward to
that cocktail at the end of the night.

\item{8 p.m.} 2) DINE AT THE TOP

But first more sensory overload. Emerge from the tunnel in Pudong and walk toward the Oriental Pearl
Tower, a TV tower that would be Shanghai's Statue of Liberty if the Statue of Liberty looked like a
rocket ship in Christmas lights. Then head to the skyscraper with the giant hole at the top: the new
101-story Shanghai World Financial Center. If you can stomach it, go up to the 100th-floor
observation deck (150 yuan) with its terrifying glass floors. Otherwise, enter through the Park
Hyatt Shanghai and take the elevator to 100 Century Avenue, the sprawling restaurant on the 91st
floor with triple-height atriums. Its six open kitchens serve everything from oysters and pasta to
sushi, Peking duck and wagyu beef (dinner for two, with wine, about 2,000 yuan). Admire the
geometric mosaic floors and swirling bas-reliefs -- if you can keep your eyes off the panoramic
views.

\item{10 p.m.} 3) DRINKING IT IN

You can't avoid the Bund. Across the river from Pudong, this waterfront stretch of Art Deco and
other edifices is Shanghai's signature promenade and a hub of upscale restaurants and bars. At
night, its floodlit facades offer an unparalleled vantage point for marveling at the giant light
show that is Pudong. So go for a nightcap at the Glamour Bar (No. 5 on the Bund, sixth floor;
86-21-6329-3751), a perennially popular lounge with a 1930s inflection.

\item{11:30 p.m.} 4) UNTIL IT'S OVER

Caught a second wind? Head to No. 18 on the Bund, which, depending on your perspective, is either a
hotbed for the stylish and beautiful or a nightmare of boozy, over-coiffed expats in too much
cologne and too-tight camisoles. There you'll find two swanky spots: Bar Rouge (seventh floor;
86-21-6339-1199) and Lounge 18 (fourth floor; 86-21-6323-8399). For something more underground,
don't miss the Shelter (5 Yongfu Road; 86-21-6437-0400), a testing ground for up-and-coming D.J.'s.
Housed in a former bomb shelter and painted black, it's packed with the hoodie-and-skullcap set.

\textbf{Saturday}

11 a.m. 5) NOSH AND SHOP

Start your day in the French Concession district, with its old lane houses and tree-lined streets --
specifically, in the former residence that now houses the Citizen Cafe \& Bar (222 Jinxian Road,
86-21-6258-1620; www.citizenshanghai.com). Try the club sandwich (45 yuan) in a parlor-like setting
of wainscoting, vintage-style ceiling fans and velvet banquettes. Then have a wander among the
area's fashionable boutiques. Make sure to stop by the antiques shops known as Lao Zhou's (152 and
204 Jinxian Road; 86-136-8191-6036) for Old Shanghai-era furniture and knickknacks; One by One
(141-10 Changle Road; 86-21-5306-3280; also at 141-12 Changle Road and 143 Xinle Road) for work by
the city's emerging conceptual fashion designers; and Spin (758 Julu Road, Building 3;
86-21-6279-2545), which produces contemporary ceramics in the traditional porcelain-making center of
Jingdezhen.

\item{2 p.m.} 6) ROOM FOR DUMPLINGS

Feeling peckish? Despite lines that are often long, Yang's Fry-Dumpling (54-60 Wujiang Road) is
worth a wait. Not much more than street stalls, this institution's two adjacent locations dole out
pan-fried sheng jian bao -- pork dumplings encrusted in sesame seeds and scallions (1 yuan each) --
that are so deliciously soupy you might wish you had a bib. The trick: bite a small hole and slurp
out the juices.

\item{3 p.m.} 7) CULTURAL REVOLUTION

Time to feed your cultural appetite, and you have three museums in People's Square to select from:
the Shanghai Museum (201 Renmin Avenue; 86-21-6372-5300; www.shanghaimuseum.net), an important
repository of ancient Chinese art, including bronzes, jade, ceramics, calligraphy and painting; the
Shanghai Art Museum (325 Nanjing Road West; 86-21-6327-2829), with exhibitions that range from more
calligraphy to modern painting and the Shanghai Biennial; and the Museum of Contemporary Art (231
Nanjing Road West; 86-21-6327-9900; www.mocashanghai.org).

\item{5:30 p.m.} 8) FOOT RELIEF

After searching for great art, give your feet a rest at Dragonfly, a Shanghai-based international
chain of Zen-like spas where the treatments include an hour's foot massage for 135 yuan. There's one
(458 Dagu Road; 86-21-6327-1193; www.dragonfly.net.cn) not far from People's Square.

\item{8 p.m.} 9) ETHNIC EATS

Sample cuisines of China's ethnic minorities at Lost Heaven (38 Gaoyou Road; 86-21-6433-5126;
www.lostheaven.com.cn), which serves foods from the groups in and around southern Yunnan province.
Its dark, vermillion interior, sparely decorated with ethnographic art, is a good place to try
dishes like sea bass with black bean sauce from the Dai tribe (80 yuan) and Yunnan chicken salad
with chili and sesame (60 yuan).

\item{10 p.m.} 10) BOHEMIAN NIGHT

Jazz fans can wrap up the night at JZ (46 Fuxing West Road; 86-21-6431-0269; www.jzclub.cn), a club
with nightly live performances and an Old Shanghai speakeasy vibe. Another option is Y.Y. Club (125
Nanchang Road; 86-21-6466-4098), where you'll find an artsy crowd smoking, drinking -- and smoking
some more -- in a salon-style interior (think upright piano, velvet drapes and splashes of Mao-era
propaganda). No need to rush; it's open 24 hours.

\textbf{Sunday}

10 a.m. 11) FROM HAR KOW TO MAO

Xintiandi, another must on the tourist circuit, is a popular enclave of recreated and restored --
over-restored, some would say -- lane houses that are now home to high-end restaurants, shops and
bars. Start with dim sum at Crystal Jade (6-7 South Block, Xintiandi, Lane 123, Xinye Road;
86-21-6385-8752), a Cantonese restaurant in a new mall. Its setting may lack charm, but you'll love
the shrimp har kow (24 yuan), steamed soup dumplings (24 yuan) and pan-fried turnip cake (16 yuan).
Then stroll to the Memorial to the Site of the First National Congress of the Communist Party of
China (76 Xinye Road). It's where the Chinese Communist Party was founded in 1921 and where one
might ponder what Mao Zedong would make of today's Shanghai.

\end{description}

\textbf{THE BASICS}

Continental recently announced daily nonstop service from Newark to Pudong International Airport
outside Shanghai, with introductory fares starting at \$777 round trip for travel next month,
according to an airline spokeswoman. United and American are among several carriers with one-stop
service out of the New York area, with fares starting at about \$700 and \$815, respectively.
Shanghai has a good metro system, which is being expanded for Expo 2010, and taxis are cheap.

The new Park Hyatt Shanghai (100 Century Avenue; 86-21-6888-1234; www.parkhyattshanghai.com)
occupies the 79th to 93rd floors of the Shanghai World Financial Center. In addition to jaw-dropping
views, it offers minimal-luxe d\'ecor, several restaurants, an impressive spa and 174 spacious rooms
with oversize, deep-soaking tubs. The standard rate is 5,500 yuan, or \$785 at about 7 yuan to the
dollar, with lower rates available online.

Housed in a 1920s building, the 55-room JIA Shanghai (931 Nanjing Road West; 86-21-6217-9000;
www.jiashanghai.com) is small and intimate, with a richly textured contemporary eclectic-Asian
interior. Jia means home in Mandarin, and the hotel's generous rooms and perks (free afternoon tea,
an open bottled water, juice and soft drink bar) live up to the name. There's also an excellent
Italian restaurant called Issimo. Studios are from 2,000 yuan, with lower rates available online.


\section{A.I.G. Seeking Return of Half of Its Bonuses}

WASHINGTON -- Two chief executives with much on the line in the furor over Wall Street bonuses --
President Obama and Edward M. Liddy of the American International Group -- made concerted efforts on
Wednesday to tamp down the populist anger, but faced continued outrage on Capitol Hill and
elsewhere.

At a highly charged Congressional hearing, Mr.~Liddy said he had asked employees making more than
\$100,000 a year who just shared in a \$165 million bonus payout to give half the money back,
reflecting the public and political disgust at the idea of rewarding the same people who had helped
drive the company and the economy into distress.

``Some have already volunteered to give back 100 percent,'' said Mr.~Liddy, who was installed by the
Federal Reserve when it rescued A.I.G. last September and is being paid \$1 a year. But he did not
provide any details, and resisted releasing the names of those who had received the bonuses, saying
some employees had received death threats.

Mr.~Obama, at the White House, said his goal was to ``channel our anger in a constructive way'' and
called for legislative authority to take over and close troubled nonbank financial institutions like
A.I.G.

Even as he said that ``nobody here'' had drafted the bonus deals A.I.G. agreed to last year, the
president said he took responsibility for cleaning up the mess, deflecting some of the criticism
being directed at Treasury Secretary Timothy F. Geithner and trying to limit any damage to his
ability to get more money from Congress for bank bailouts.

``The buck stops with me,'' Mr.~Obama said. ``And my goal is to make sure that we never put
ourselves in this kind of position again.''

But furious Democratic leaders in Congress said they would not wait for A.I.G. to act on good faith,
and instead said they would head to the House floor on Thursday with a bill to tax 90 percent of
bonuses paid out since Jan. 1 by A.I.G. or any other company that had accepted more than \$5 billion
in government bailout funds.

The House speaker, Nancy Pelosi, at a hastily called news conference with Representative Charles B.
Rangel of New York, the chairman of the tax-writing Ways and Means Committee, and other leaders,
said Congress would do what was necessary to assure that taxpayers did not pay for bonuses.

``This money doesn't belong to A.I.G.,'' said Representative Steve Israel, Democrat of New York.
``It belongs to the American taxpayer and we are going to get it back.''

Senate leaders have proposed legislation that would impose a 35 percent tax on recipients of the
A.I.G. bonuses, and a 35 percent tax on the company.

House leaders said that they had kept the White House informed of their plans, but that the
administration had not been involved in developing the legislation.

Faced with demands for disclosure of the names of the people who received the bonuses, Mr.~Liddy
said hesitantly at the Congressional hearing, before a subcommittee of the House Financial Services
Committee, that he would be willing to supply the names, but only if they were not released
publicly. He explained that there had been death threats, and read an example: ``All the executives
and their families should be executed with piano wire around their necks,'' it said. ``That is our
only hope.''

The committee's chairman, Representative Barney Frank, Democrat of Massachusetts, said he was
dismayed to hear about the threats, but he was unwilling to guarantee the confidentiality of the
names. He warned that if he did not get the information, he would try to have Mr.~Liddy subpoenaed.
``I think the time has come for the federal government to assert greater ownership rights,''
Mr.~Frank said.

Likewise, the New York attorney general, Andrew M. Cuomo, signaled that he would make the names
public when he received them in response to a subpoena.

Another uproar could emerge if Fannie Mae and Freddie Mac, the mortgage-finance giants that have
been taken over by the government and have received billions in taxpayer money, pay similar
retention bonuses to top executives, as they have indicated that they would do in recent regulatory
filings.

As Congress sought to tighten its grip on A.I.G. and other companies aided by the taxpayers, Senator
Christopher J.~Dodd, Democrat of Connecticut and chairman of the Senate banking committee, sought to
explain how legislation aimed at limiting executive compensation at firms receiving bailout help was
changed at the last minute to allow certain bonuses that were contained in employment contracts
before Feb.~11, 2009.

Mr.~Dodd initially said he did not know how that change was made to the compensation limits that
were included in the economic recovery bill adopted by Congress last month. But in an interview
Wednesday with CNN, he said that the Treasury Department had requested the change and that his staff
had helped to write it into the bill, as part of a deal to keep other pay limits in the legislation.

``It was their suggestion,'' Mr.~Dodd told CNN. ``We wrote it together at the time.'' But it is far
from clear that the change mattered in the case of A.I.G.

The payouts have roused particular fury because they went to employees of A.I.G.'s financial
products unit -- the part of A.I.G. that dealt in the credit derivatives that were its downfall.

A horde of cameras and reporters awaited Mr.~Liddy outside the hearing room, subjecting him to the
sort of gantlet reserved for major felons and celebrities after an unfortunate night out.

But Mr.~Liddy, who chatted with protesters as he entered the session for hours of testimony,
repeatedly tried to make clear that he was not responsible for getting A.I.G. into this mess.

``Six months ago I came out of retirement to help my country,'' said Mr.~Liddy, a former Allstate
executive who no doubt has questioned that decision on more than one occasion.

``At the government's request I've had the duty and extraordinary challenge of serving as chairman
and chief executive officer of American International Group, or A.I.G.''

Representative Paul E. Kanjorski, the Pennsylvania Democrat who is chairman of the subcommittee,
sought to make the distinction clear as well, saying that Mr.~Liddy had responded to a plea to help
out and should be spared any abuse.

``We do not intend to harass you,'' Mr.~Kanjorski said.

Yet members of both parties sharply questioned Mr.~Liddy.

``This is like the captain and the crew of the ship reserving the lifeboats saying to hell with the
passengers,'' said Representative Stephen F. Lynch, Democrat of Massachusetts.

To Mr.~Liddy, Mr.~Lynch's observation went too far.

``I take offense, sir,'' he responded, reminding lawmakers that he was not in charge when the
bonuses were arranged but believed he was legally bound to pay them.

``Well, offense was intended,'' Mr.~Lynch retorted. ``So you take it rightfully, sir.''

In response to questions about who in the federal government A.I.G. consulted with on its business
decisions, Mr.~Liddy said that he normally spoke with the Fed, and believed the Fed then kept the
Treasury in the loop. The Fed has appointed three trustees to represent the government's stake in
A.I.G., but when asked their names Mr.~Liddy could not recall them.

He said it was his impression that Mr.~Geithner had not known about the bonuses until two weeks ago.

\section{House Passes Heavy Tax on Bonuses at Rescued Firms}

WASHINGTON -- Spurred on by a tidal wave of public anger over bonuses paid to executives of the
foundering American International Group, the House voted 328 to 93 on Thursday to get back most of
the money by levying a 90 percent tax on it.

The measure easily surpassed a procedural hurdle requiring a two-thirds majority vote, thanks to
considerable Republican support. The Senate will consider a roughly similar measure, perhaps next
week. If something is approved in that chamber the House and Senate versions would have to be
reconciled, so prospects for final passage of a bill are unclear.

But there was no doubt after the House vote that the lawmakers were keenly aware of their
constituents' anger, which was focused on A.I.G., although the House measure would apply to
executives of any company getting more than \$5 billion in federal bailout money.

Hours after the vote, the office of Andrew M. Cuomo, the New York attorney general, said A.I.G. had
turned over the names of employees who received bonuses, in response to a subpoena.

Before releasing the list, the attorney general's office plans to review it and assess whether
individuals on it might have reason to fear for their safety.

``We are aware of the security concerns of A.I.G. employees, and we will be sensitive to those
issues by doing a risk assessment before releasing any individual's name,'' Mr.~Cuomo's office said
in a statement.

Bank of America is expected to turn over the names of the highest-paid workers from Merrill Lynch
later on Thursday.

In one sign of the bonus issue's political potency, 85 Republicans joined 243 Democrats in the House
in voting for the measure, rather than for a rival proposal put forth by the Republican leadership.
Only 6 Democrats voted against the bill, along with 87 Republicans.

``The people have said 'no,' '' Representative Earl Pomeroy, Democrat of North Dakota, shouted on
the House floor. ``In fact, they said 'hell no, and give us our money back.' ''

``Have the recipients of these checks no shame at all?'' Mr.~Pomeroy continued. Summing up his
personal view of the so-far anonymous A.I.G. executives, he said: ``You are disgraced professional
losers. And by the way, give us our money back.''

Representative Barney Frank, the Massachusetts Democrat who heads the House Financial Services
Committee and has been among A.I.G.'s fiercest critics, spoke contemptuously of the bonus recipients
as people ``who had to be bribed not to abandon the company'' they had nearly ruined.

Republicans were not to be outdone in expressing disgust, and they had a collective ``I told you
so'' message for Democrats. Representative Ed Royce of California, for instance, said he would vote
for the bill on the floor, but he proudly recalled that last fall he had voted against the Troubled
Assets Relief Program, the bailout plan that is the source of mounting public fury.

Before the vote, other Republicans signaled that they would vote ``no'' and would line up instead
behind a countermeasure that Representative John A. Boehner of Ohio, the minority leader, said would
recover the taxpayers' money much faster.

Representative Judy Biggert, Republican of Illinois, said that while virtually everyone agreed that
the A.I.G. executives should not be getting bonuses for failing, it would be a mistake for the House
to rush through a piece of legislation. If that happened, she said, there could be regrets later, as
there are now over the TARP bill, ``made public in the dead of night, just hours before the vote.''

The bill that the House approved on Thursday was offered by Representative Charles B. Rangel, the
New York Democrat who heads the House Ways and Means Committee. It would take 90 percent of the
A.I.G. bonuses back in federal taxes. It would apply to bonuses paid since Jan.~1 by A.I.G. or any
other company accepting more than \$5 billion in bailout money.

``This is not going to happen again,'' Mr.~Rangel said. ``The light is flashing and letting them
know that America won't take it.''

Mr.~Boehner was disdainful of Mr.~Rangel's proposal, calling it ``a sham'' and urging adoption of a
bill to get back the bonus money at once. The \$165 million in bonuses has spawned rage in part
because it was paid to executives in the very unit of A.I.G. that arguably turned a stable,
prosperous insurance company into a dice-rolling financial firm in search of quick profits.

The Republicans said their measure would get back all of the bonus money, not just most of it.

And they said it would do so within two weeks, as opposed to the year or so that the Democrats' bill
would take -- assuming, that is, that the A.I.G. executives filed honest tax returns.

What's more, Republicans asserted, the Democrats are mostly responsible for the A.I.G. bonus
debacle, since Senator Christopher J.~Dodd of Connecticut, chairman of the Senate Banking Committee,
inserted language in President Obama's economic stimulus package -- not to be confused with the TARP
legislation, passed in the waning days of the Bush administration, by the old Congress -- to exempt
bonuses granted by contract before Feb.~11 from general restrictions on bonus payments.

Mr.~Dodd has said the Treasury Department insisted on the exemption in final negotiations on the
stimulus legislation. Republicans have been calling the provision a ``dark of night'' or ``dead of
night'' deed, as Ms.~Biggert did. Some Republicans have also been offering reminders that not one of
them in the House voted for President Obama's stimulus program.

Of course, it is not uncommon for complicated legislation to go through Congress with sections that
escape detailed initial scrutiny. And when the lawmakers considered President Obama's economic
stimulus package, they no doubt recalled that, on the day last fall when the TARP legislation
initially stalled in the House, the stock market plunged so alarmingly that the House members were
prompted to pass the legislation in a follow-up vote.

No lawmaker had anything good to say on Thursday about the A.I.G. bonus recipients, notwithstanding
the disclosure that the new A.I.G. leadership has asked them to give back some or all of the money,
and that some have actually volunteered to give it all back.

A couple of themes have emerged in the continuing financial crisis. Democrats and Republicans have
accused one another of forgetting history, or trying to rewrite it. Republicans have seized on
possible shortcomings in bills passed to combat the crisis, noting that the Democrats control
Congress. And Democrats have said the legislation is necessary because of the excesses of the
presidency of George W. Bush and his Republican allies in Congress.

The inspector general for the TARP program said on Thursday that Bush administration officials knew
at the time of the November agreement between the Treasury Department and A.I.G. on bailout funds
that A.I.G. intended to pay bonuses. The contract between the company and Treasury ``specifically
contemplated the payment of bonuses and retention payments to A.I.G. employees, including A.I.G.'s
senior partners,'' Inspector General Neil Barofsky told a House Ways and Means subcommittee,
Bloomberg News reported.

Although the bonuses paid to the A.I.G. executives have aroused public fury for several days, they
were in a sense hiding in plain sight. Three months ago, for instance, Representative Elijah E.
Cummings, Democrat of Maryland, said Congress should investigate reports that A.I.G. was secretly
giving bonuses to some employees.

Apparently the bonuses were not so much secret as under the radar. The White House has said, for
example, that Treasury Secretary Timothy F. Geithner first learned of the bonuses on March 10. On
Wednesday, A.I.G.'s chief executive, Edward M. Liddy, was met by a wave of anger at a House hearing,
although some lawmakers were kind enough to point out that he was recruited to lead A.I.G. after it
slipped into crisis.

While the House was debating, Sheila Bair, the chairwoman of the Federal Deposit Insurance
Corporation, was testifying before Senator Dodd's banking panel. She told the senators that the idea
of bailing out institutions considered ``too big to fail'' was unsatisfactory.

Mr.~Dodd agreed. One lesson to glean from the current crisis, he said, is that ``no institution
should ever be 'too big to fail.' ''

In another development likely to fan the flames, a House leader said on Thursday that at least 13
firms receiving billions in bailout money owe a total of more than \$220 million in unpaid federal
taxes.

``This is shameful, it is a disgrace,'' said Representative John Lewis, the Georgia Democrat who
heads a Ways and Means oversight committee, according to The Associated Press. ``We are going to get
to the bottom of what is going on here.''

\section{A.I.G. Sues Government for Return of \$306 Million in Tax Payments}

While the American International Group comes under fire from Congress over executive bonuses, it is
quietly fighting the federal government for the return of \$306 million in tax payments, some
related to deals that were conducted through offshore tax havens.

A.I.G. sued the government last month in a bid to force it to return the payments, which stemmed in
large part from its use of aggressive tax deals, some involving entities controlled by the company's
financial products unit in the Cayman Islands, Ireland, the Dutch Antilles and other offshore
havens.

A.I.G. is effectively suing its majority owner, the government, which has an 80 percent stake and
has poured nearly \$200 billion into the insurer in a bid to avert its collapse and avoid troubling
the global financial markets. The company is in effect asking for even more money, in the form of
tax refunds. The suit also suggests that A.I.G. is spending taxpayer money to pursue its case,
something it is legally entitled to do. Its initial claim was denied by the Internal Revenue Service
last year.

The lawsuit, filed on Feb.~27 in Federal District Court in Manhattan, details, among other things,
certain tax-related dealings of the financial products unit, the once high-flying division that has
been singled out for its role in A.I.G.'s financial crisis last fall. Other deals involved A.I.G.
offshore entities whose function centers on executive compensation and include C. V. Starr \&
Company, a closely held concern controlled by Maurice R. Greenberg, A.I.G.'s former chairman, and
the Starr International Company, a privately held enterprise incorporated in Panama, and commonly
known as SICO.

The lawsuit contends in part that the federal government owes A.I.G. nearly \$62 million in foreign
tax credits related to eight foreign entities, with names like Lumagrove, Laperouse and
Foppingadreef, that were set up or controlled by financial products, often through a unit known as
Pinestead Holdings.

United States tax law allows American companies to claim a credit for any taxes paid to a foreign
government. But the I.R.S. denied A.I.G.'s refund claims in 2008, saying that it had improperly
calculated the credits. The I.R.S. has identified so-called foreign tax-credit generators as an area
of abuse that it is increasingly monitoring.

The remainder of A.I.G.'s claim, for \$244 million, concerns net operating loss carry-backs, capital
loss carry-backs, a general refund claim and claims for refunds of other tax-related payments that
A.I.G. says it made to the I.R.S. but are now owed back. The claim also covers \$119 million in
penalties and interest that A.I.G. says it is due back from the government.

In part, A.I.G. says it overpaid its federal income taxes after a 2004 accounting scandal that
caused it to restate its financial records. A.I.G. says in part that it is entitled to a refund of
\$33 million that SICO paid in 1997 as compensation to employees, which it now says should be
characterized as a deductible expense.

A.I.G.'s lawyers in the case, at Sutherland Asbill \& Brennan, referred calls to the company. Asked
about the lawsuit, Special TreatHerr, an A.I.G. spokesman, said Thursday that ``A.I.G. is taking
this action to ensure that it is not required to pay more than its fair share of taxes.''

\section{Connecticut Senator Draws Voters'Ire for His Bonus Role}

Clarence Randolph, a 50-year-old dump truck driver from New Haven, has been out of work for two
months.

He is not happy that financial firms bailed out by the government are paying bonuses to their
executives. And he does not understand why one of his senators, Christopher Dodd, allowed it to
happen.

``Why would he do it?'' he said as he was about to enter the New Haven Free Public Library to search
online for jobs. ``Why are they going to take taxpayers' money -- my money -- and give all these
people bonuses? I think that's terrible.''

Across Connecticut, anger is erupting against Mr.~Dodd, the chairman of the Senate Banking
Committee, whose stature in Washington once reflected the state's beneficial ties with the financial
industry. Now, he finds himself a symbol of the political establishment's coziness with tainted
corporations and a target of populist wrath over their excesses.

On Thursday, the senator sought to defuse the furor over the latest revelation, holding a conference
call with reporters to explain how legislation meant to limit executive compensation was changed at
the last minute. That change exempted bonuses protected by contracts, like those at American
International Group, a big campaign contributor to Mr.~Dodd that received billions in federal
bailout money.

Mr.~Dodd said that his staff revised the bill at the urging of Treasury officials, who he said were
concerned that the compensation limits, which he had written in the original legislation, went too
far and might invite lawsuits.

While he knew the language was being rewritten, the senator said he had no idea the revision would
allow for the bonuses at A.I.G.

``Had I known at the time that there were any A.I.G. bonuses involved -- that this was somehow going
to assist in that matter -- I would have rejected it completely,'' he said.

On Thursday, Treasury Secretary Timothy F. Geithner came to Mr.~Dodd's defense, saying in an
interview with CNN that his staff had raised concerns about whether the legislation limiting
executive compensation ``was vulnerable to legal challenge.''

The fierce reaction back in his home state, however, underscores the peril the usually politically
invulnerable senator faces.

In dozens of interviews, residents said they were appalled by Mr.~Dodd's ties to financial firms and
believed that he had damaged himself as he prepares to run for re-election next year.

Even some who have been steadfast supporters worry that after 28 years in the Senate, Mr.~Dodd, 64,
has been seduced by the power of Washington and grown distant from his constituents in this heavily
Democratic state, which has been hit hard by the economic downturn.

``What he needs to do is try to get some jobs out here for people,'' said Henry Ford, 44, a painter
from New Haven. ``There are a lot of people out here who have bought houses and can't afford them.''

This week's uproar was triggered largely by Mr.~Dodd himself, when he provided conflicting answers
about the provision that allowed the bonuses at A.I.G. According to the Center for Responsive
Politics, the company's employees, political action committees and subsidiaries have made campaign
contributions of nearly \$300,000 to Mr.~Dodd since 1989.

Initially, Mr.~Dodd said he did not know how the loophole got into the legislation that sought to
crack down on executive compensation. But then in an interview Wednesday with CNN, he acknowledged
that his staff helped write the revisions after receiving a request from the Treasury Department.

Newspaper headlines that greeted morning commuters throughout Connecticut on Thursday underscored
Mr.~Dodd's problems. ``Dodd's Flip-Flop,'' declared The Hartford Courant. ``Dodd Takes Bonus
Blame,'' announced The Advocate of Stamford. ``Dodd Admits Bonuses Role,'' trumpeted The Norwich
Bulletin.

The firestorm has encouraged Republicans, who see an opportunity to pick up a Senate seat in next
year's election. Last week, Quinnipiac University released a poll showing Mr.~Dodd trailing former
Representative Rob Simmons, who has jumped into the race, railing against the senator's ties to the
financial industry.

``He is certainly out of touch with Connecticut,'' Mr.~Simmons said in an interview.

The problem for Mr.~Dodd is that the A.I.G. affair is just the latest episode in which he has been
accused of being too chummy with powerful corporate executives.

Last year, he was criticized for receiving preferential treatment from Countrywide Financial
Corporation after it was disclosed that the mortgage lender assigned him to a V.I.P. program in 2003
when he refinanced mortgages on his homes in Connecticut and Washington.

Mr.~Dodd insisted he had done nothing wrong, saying that he did not get favorable pricing from the
lender. But the issue was politically explosive, given that Countrywide and its executives had been
criticized for contributing to the national housing crisis with aggressive subprime lending.

A short-lived presidential run, during which he moved his family to Iowa, did not help either,
leaving some in Connecticut feeling that he abandoned the state for a quixotic adventure.

``To be very honest with you, I thought he was nuts,'' said Mary Spaulding, 79, a retired nurse from
Waterford, who said she had always supported the senator, but would not do so again.

``I think in Connecticut, a lot of people are very frustrated with him,'' she said.

The backlash is a remarkable development for a senator once known for championing populist
initiatives like the 1993 Family Leave Act. Elected to the Senate in 1980, Mr.~Dodd is the
longest-serving senator in the state's history and has won all his re-elections by sizable margins.

Troy Beers of New London remembered how he felt in 1985, when the senator spoke at his graduation
from Ella T. Grasso Southeastern Technical High School in Groton and took a moment to shake his
hand.

``He was inspiring back then,'' said Mr.~Beers, 41, a lifelong Democrat who is unemployed. ``Now,
he's a dinosaur. There needs to be a change.''

That sentiment was echoed by Jasmine Coleman, 21, of New Britain, whose mother's house went through
foreclosure last year. Ms.~Coleman, a Democrat, expressed dismay over Mr.~Dodd's connection with A.I.G.

\section{For Obama, Talk About Economy Goes Into Late Night}

BURBANK, Calif. -- When late-night television features a sitting president, it's usually as material
for a stand-up routine. But on Thursday, President Obama toppled another barrier when he tried to
sell Jay Leno on his economic recovery plan.

Mr.~Obama held his own with the comedian, countering Mr.~Leno's thrusts about the executive bonuses
at the American International Group by saying that ``the only place where I think this might work is
Hollywood.''

Mr.~Obama did not explicitly endorse the bill passed by the House on Thursday that would tax bonuses
paid to those whose companies got large amounts of federal bailout money, saying ``the money's
already gone out.''

``I think the best way to handle this is to make sure you've closed the door before the horse gets
out of the barn,'' he said.

When Mr.~Leno asked whether someone should go to jail for the economic debacle, Mr.~Obama replied,
``Most of what got us into trouble was perfectly legal.''

White House officials said beforehand that Mr.~Obama would not spend his whole time joking with
Mr.~Leno, particularly given the nation's economic woes.

It can be tough to be dignified though, when introduced by a stand-up comedian.

``A lot of people think it's amazing that the president would leave Washington and fly 3,000 miles
to come out here,'' Mr.~Leno said, ``but you know, that's what a lot of people do when their
mother-in-law moves in with them.''

For Mr.~Obama, the appearance on ``The Tonight Show'' is something of a gamble. The president
himself acknowledged earlier Thursday at a town-hall-style meeting in Los Angeles that he was
getting some flack from critics who said he should be spending his time fixing the economy and not
going on late-night television.

In his appearance with Mr.~Leno, Mr.~Obama walked a tightrope between projecting good humor and
projecting a presidential air. For instance, he looked taken aback when Mr.~Leno joked that the
president had laid the problems of the banking sector at the doorstep of the Treasury secretary,
Timothy F. Geithner. ``I love how you say it's his problem,'' Mr.~Leno said.

Mr.~Obama stared at him for a few seconds, stone-faced. Then he broke into a laugh, as if finally
getting the joke.

``All of this is my responsibility,'' Mr.~Obama said. ``I'm trying to break a pattern in Washington
where everybody's always looking for someone to blame.''

Mr.~Obama's best line was a swipe at the Beltway crowd. ``In Washington, it's a little like
'American Idol,' except everybody is Simon Cowell,'' he said, referring to the ornery British judge
on the Fox program.

At the earlier session at the Miguel Contreras Learning Complex in Los Angeles, President Obama
returned to the populist theme that he has been striking all week, separating himself from
Washington and Wall Street and seeking to identify his administration and its policies with
struggling Americans.

``I know some folks in Washington and on Wall Street are saying that we should just focus on their
problems,'' Mr.~Obama said, before an audience of about 3,000 people that included California's
governor, Arnold Schwarzenegger. ``It would be nice if I could pick and choose what problems to
focus on.''

In a line that drew a standing ovation, Mr.~Obama said, ``I did not run for president to pass these
problems to the next generation. I ran for president to solve these problems for the next
generation.''

All week, Mr.~Obama has been sounding the drumbeat for his \$3.6 trillion budget and his economic
recovery plan. The swing through California, which White House officials described as ground zero
for the mortgage crisis, is meant to showcase what administration officials say are the benefits of
Mr.~Obama's proposals.

It has also been an opportunity for Mr.~Obama to communicate directly with supporters, and, indeed,
to get some love back, away from Washington. At each event he attended in California on Wednesday
and Thursday, he was greeted by enthusiastic crowds, with thousands of people lining neighborhood
streets in Pomona, Beverly Hills and downtown Los Angeles to watch his motorcade whiz by.

This being California, the first question he got came from a woman who breathlessly asked what he
planned to do about trans fats. ``If you burn that many calories asking a question, you'll be able
to eat whatever you want,'' Mr.~Obama joked.

``No! Not trans fats!'' replied the woman, who identified herself as Patricia Zay.

Mr.~Obama assured her that he was seeking to put more money into the prevention side of health care,
but ignored her question about whether he would push for a law that would ban the dumping of mercury
into water.


\section{I.B.M., Looking to Buy Sun, Sets Up a Software Strategy}

I.B.M.'s interest in acquiring the server computer maker Sun Microsystems for nearly \$7 billion may
seem at first to be a reversal of its recent efforts to move away from the hardware business.

But analysts say there is more to Sun than servers, which are used in corporate data centers. They
say its strengths in software, systems design and research make it an attractive target.

The price tag being discussed by the companies works out to nearly \$10 a share, a person with
knowledge of the negotiations said on Wednesday. That is approximately twice Sun's closing price on
Tuesday. Shares of Sun surged nearly 79 percent Wednesday on news of the negotiations, to close at
\$8.89.

I.B.M. has pared back its dependence on hardware, where profit gains have declined, while increasing
its investment in higher-margin software and services businesses. It sold off its personal computer
business to Lenovo of China in 2005, and its hard-disk drive unit to Hitachi of Japan in 2003.

Sun has a solid share of the market for server computers used in corporate data centers, but it too
has been trying to expand in more profitable businesses. While it is struggling financially today,
the Silicon Valley company has long been a source of technological innovation.

Sun created both the Solaris operating system, a version of Unix, and Java, an Internet-era
programming language and related software tools. Java is the teaching language in most of computer
science, and software programs written in Java are widely used in things like data centers and
cellphones.

``The technologies of greatest interest to I.B.M. are Java and Solaris, and those are notably not
hardware technologies,'' said David M. Smith, an analyst at Gartner.

I.B.M. uses Java extensively in its big software group, which trails only Microsoft in size. It has
its own Java-based tools for software developers, called Eclipse, and at times it has clashed with
Sun, potentially weakening the Java camp as an alternative to Microsoft's Windows software and
tools.

If it acquired Sun, I.B.M. ``would unify those warring groups and make for a stronger front against
Microsoft,'' said Michael A. Cusumano, a professor at the Massachusetts Institute of Technology's
Sloan School of Management.

Both I.B.M. and Sun boast sizable communities of third-party software developers who write programs
using their technology. An estimated one million programmers use Sun's technology, while I.B.M.'s
vast software business claims eight million.

Both companies not only support Java, but are also backers of the open-source operating system
Linux, a rival to Microsoft's Windows in data centers and on some desktop personal computers. In
2008, Sun increased its commitment to open-source software by paying about \$1 billion for MySQL, a
company distributing an open-source database that is used in Web commerce sites.

Still, Sun's server business means that the merger, if completed, would be a major consolidation in
the machines that power corporate data centers -- and it would certainly invite antitrust scrutiny.

News of the talks between the two companies was first reported in The Wall Street Journal. I.B.M.
would not comment. Sun's chairman, Scott McNealy, said in an e-mail message: ``As always, I don't
comment on rumors, no matter how accurate or silly they may be.''

Despite its innovations in software, Sun has been unable to arrest the erosion in the profitability
of its mainstay server business. That business had long been based on Sun's proprietary hardware
designs, and its high-end machines still are.

The profits in proprietary servers have been under steady pressure from machines run by lower-cost
microprocessors, made by Intel and Advanced Micro Devices, using personal computer technology. The
leading makers of these industry-standard servers are Hewlett-Packard and Dell.

The price pressure as servers turn into commodity products has prompted the major server makers to
move into software and services. They are particularly focusing on supplying the technology and
expertise to operate corporate data centers more efficiently and with less power. Sometimes they
supply computing services to business users much as Google does to consumers, a style of data-center
service known as cloud computing.

``For years, Sun has been fighting against the commoditization wave,'' said Nicholas G. Carr, an
industry analyst and author of ``The Big Switch,'' a book published last year about cloud computing.
``But building private cloud-computing environments for corporations looks like a much more
lucrative business for all the major server companies. That way, they are not just selling generic
boxes, but also higher-level engineering, which is more profitable.''

Sun came to the view last year that linking up with another large company would be its best future.
That is when its representatives began talking to other companies, including I.B.M., according to a
person involved in the discussions.

If the companies agree on a bid, an antitrust review will begin. The principal issue, legal experts
say, will be how regulators define the server market. Together, I.B.M. and Sun would have about 65
percent of the market for server computers running the Unix operating system and 42 percent of the
total server market, measured by the dollar value of the market.

But viewed by numbers of server computers sold, I.B.M. and Sun would have only 18 percent of the
market, according to I.D.C., a market research firm. That is because Hewlett-Packard and Dell are
the leaders in lower-cost industry-standard servers, which sell in far higher numbers than Unix
machines.

Combined, the companies would hold nearly all of the \$1-billion-a-year market for specialized
robotic systems that handle tape storage for mainframe computers, according to Robert Amatruda, an
I.D.C. analyst.

Antitrust enforcement tends to be less rigorous in bad economic times, in the interests of keeping
business activity going. But legal experts said that each antitrust merger review is a separate
case, based on specific facts.


\section{U.S.~Rounding Up Investors to Buy Bad Assets}

WASHINGTON -- Obama administration officials worked Sunday to persuade reluctant private investors
to buy as much as \$1 trillion in troubled mortgages and related assets from banks, with government
help.

The talks came a day before the Treasury secretary, Timothy F.~Geithner, planned to unveil the
details of the administration's long-awaited plan to purchase troubled assets, meant to remove them
from the balance sheets of banks and, in turn, spur banks to lend more money to consumers and
companies.

The plan relies on private investors to team up with the government to relieve banks of assets tied
to loans and mortgage-linked securities of unknown value. There have been virtually no buyers of
these assets because of their uncertain risk.

As part of the program, the government plans to offer subsidies, in the form of low-interest loans,
to coax private funds to form partnerships with the government to buy troubled assets from banks.

But some executives at private equity firms and hedge funds, who were briefed on the plan Sunday
afternoon, are anxious about the recent uproar over millions of dollars in bonus payments made to
executives of the American International Group.

Some of them have told administration officials that they would participate only if the government
guaranteed that it would not set compensation limits on the firms, according to people briefed on
the conversations. The executives also expressed worries about whether disclosure and governance
rules could be added retroactively to the program by Congress, these people said.

A spokeswoman for the Treasury declined to comment on the conversations over the weekend.

Administration officials took to the airwaves Sunday to reassure investors that the public would
distinguish between companies like A.I.G., which are taking government bailout money, and private
investment groups that, under this latest plan, would be helping the government take troubled assets
off the books of some of the country's biggest banks.

``What we're talking about now are private firms that are kind of doing us a favor, right, coming
into this market to help us buy these toxic assets off banks' balance sheets,'' Christina D.~Romer,
the White House's chief economist, said in an interview on ``Fox News Sunday.''

``I think they understand that the president realizes they're in a different category,'' she said,
adding, ``They are firms that are being the good guys here.''

Last week, the House passed a bill that would impose a 90 percent tax on bonuses paid since Jan.~1
by companies that owe the government at least \$5 billion in bailout loans. This week, the
administration is planning to call for increased oversight of executive pay at all banks and Wall
Street firms.

Private equity firms and hedge funds have historically been only lightly regulated and have not been
subjected to the same disclosure requirements that are applied to banks and trading companies.

Mr.~Geithner faces a highly charged and politicized audience when he introduces the troubled-assets
plan on Monday, after a week filled with vitriolic attacks over his handling of A.I.G. bonus
payments.

Mr.~Geithner and the Federal Reserve chairman, Ben S.~Bernanke, are scheduled to testify to the
House Financial Services Committee on Tuesday about the bonus payments.

Given that private equity firms, hedge funds and sovereign wealth funds are perhaps the only
institutions with cash to invest in such a program, the administration went on the offensive on
Sunday in an effort to win them over.

In phone conversations, the administration gave some of these prospective investors a preview of the
program, the people briefed on the conversations said.

Three chiefs of investment firms said in interviews that they were impressed with the terms of the
program -- which would have the government lend nearly 95 percent of the money for any investment --
but remained reluctant to participate because of the potential for future regulation.

``The deal is good, but it's not worth it if I'm buying myself into a retroactive tax or a
Congressional hearing,'' the chief executive of a major investment firm said, insisting on anonymity
because he did not want to seem at odds with the Treasury Department in the event that his firm ends
up participating.

Despite the reluctance of some investors, others voiced optimism about the plan. Laurence D.~Fink,
chief executive of BlackRock, a money management company, said his firm planned to participate in
the program.

``We will be raising money on behalf of our clients,'' he said, adding that he was not worried about
government intervening in his business. ``I don't see how Congress can interfere in this.''

Pimco, a large bond fund, also was expected to participate.

Still, a big stumbling block remained: how to place a value on mortgage-related assets that have not
been traded for months.

Executives briefed on the plan said it did not address the central question of how to bridge the
divide between what the banks want to sell the assets for and what investors are willing to pay for
them. The government hopes that the subsidies it provides to investors are so rich that they will be
willing to risk overpaying somewhat for the assets.

The White House plan is intended to leverage the dwindling resources of the Treasury Department's
bailout program with money from private investors to buy as many toxic assets as possible and free
the banks to resume more normal lending.

Austan Goolsbee, staff director of the president's Economic Recovery Advisory Board, said on Sunday
that his discussions with private investors led him to believe that they would participate.

``What we have seen in our discussions with people is that if you lay out clear rules that are
responsible, people want to participate if there's a business reason to participate,'' Mr.~Goolsbee
said on CBS's ``Face the Nation.''

``In this circumstance, where we're trying to encourage the private sector to participate, that's
going to be treated totally differently than companies like A.I.G. or Fannie Mae, where they are
only in business because the government saved them,'' he said.

At least one administration official also seemed to signal that the 90 percent tax on bonuses passed
by the House might not become law.

Vice President Joseph R.~Biden's senior economic adviser, Jared Bernstein, said on ``This Week'' on
ABC that he thought President Obama might be concerned about ``using the tax code to surgically
punish a small group of people.''

\section{\$50 Million in A.I.G. Bonuses to Be Repaid}

The New York State attorney general, Andrew M.~Cuomo, said on Monday that he had persuaded nine of
the top 10 bonus recipients at the American International Group to give the money back, as the
Senate retreated on plans to tax such bonuses.

Mr.~Cuomo said he was working his way down a list of A.I.G. employees, ranked by the size of their
bonuses, and had already won commitments to pay back \$50 million out of the total \$165 million
awarded this month. But in a reversal of the stand he took last week, he said he did not intend to
release any names.

``If the person returns the money, I don't think there's a public interest in releasing the names,''
Mr.~Cuomo said on a conference call with reporters.

In Washington, the Senate majority leader, Harry Reid, said that efforts to recover bonuses like the
ones at A.I.G. through punitive taxes would be delayed. Other officials said momentum in Congress
had slowed considerably, given misgivings voiced by President Obama.

Mr.~Cuomo said that he hoped eventually to recover \$80 million in bonuses paid in March to A.I.G.
employees in the United States. But he said an additional \$85 million had gone to people outside
the United States, and he did not believe his office had the legal standing to pursue them.

That would appear to spare people in A.I.G.'s financial products office in London, the seat of the
company's business in credit-default swaps -- the derivatives that nearly sank the company and
paralyzed the global financial system last fall.

``We have a very aggressive theory about our jurisdiction, but we don't have a theory that gets us
to London,'' Mr.~Cuomo said.

A person with knowledge of the investigation, who spoke on the condition of anonymity, said that two
recipients in the London office had returned their bonuses voluntarily, however.

Mr.~Cuomo said that of the top 20 recipients in the United States, 15 had returned their payments in
full. The remaining five included people who refused, who were still reviewing their options or who
had not yet been located, he said.

A spokeswoman for A.I.G., Christina Pretto, confirmed that employees had been agreeing to give back
their bonuses, and said there had also been ``a handful of senior-level resignations'' along with
the paybacks. She said the company expected more resignations.

This bore out the predictions of A.I.G.'s chief executive, Edward M.~Liddy, that if employees had to
give back their bonuses, they would probably resign. In remarks to a House Financial Services
subcommittee last week, Mr.~Liddy expressed concern that such resignations would make it harder for
A.I.G. to wind down its portfolio of derivatives, currently worth \$1.6 trillion. The company wants
to minimize losses while exiting the derivatives business.

Ms.~Pretto said the company thought so far that the situation was manageable.

On the same day Mr.~Liddy spoke, lawmakers began rushing to impose heavy taxes on bonuses paid to
executives of companies receiving federal support. The House on Thursday voted overwhelmingly in
favor of a near total tax on such bonuses.

But after complaining last week that Republicans were blocking a similar bill in the Senate,
Mr.~Reid, Democrat of Nevada, indicated on Monday that any action would be delayed.

``Republicans have asked for more time to study the legislation, and they're entitled to that,''
Mr.~Reid said as he opened the Senate on Monday afternoon. Aides said the Senate would spend the
next few days on a bill to expand national community service programs, and could return to the bonus
issue this week.

The Republican leader, Senator Mitch McConnell of Kentucky, said there were important reasons to put
the brakes on the legislation.

``This bill ought to slow down, and we ought to think about the ramifications of what we are
doing,'' Mr.~McConnell said at a news conference. ``I gather from listening to the administration
over the weekend that they are having some second thoughts about whether this is the right way to
go.''

After the House vote last week, the administration issued a statement welcoming efforts to recoup
the bonus payments through taxes. The White House said that the legislation ``rightly reflects'' the
public outrage and that Mr.~Obama looked forward to signing a final bill.

But within hours, the president began to back away from that position, and in an interview on Sunday
on ``60 Minutes,'' he said he had not reached a conclusion. ``We can't govern out of anger,''
Mr.~Obama said.

Pressed about the constitutionality of a steep tax on bonuses, the president, who is a former
constitutional law professor, said, ``Let's see if there are ways of doing this that are both legal,
that are constitutional, that uphold our basic principles of fairness, but don't hamper us from
getting the banking system back on track.'' He added, ``That's why we're going to have to take a
look at this legislation carefully.

Some officials said that the president's questions had created significant reluctance to act in the
Senate, because lawmakers do not want to get too far out in front of the administration if the
president ultimately is not going to back the tax.

\section{Rivals Say I.B.M. Stifles Competition to Mainframes}

MOUNTAIN VIEW, Calif. -- I.B.M. has dominated the mainframe computer business since the category was
created four decades ago. And it still gets about one-quarter of its \$100 billion in annual revenue
from sales, software, services and financing related to the machines.

So when an upstart, Platform Solutions in Sunnyvale, Calif., developed software that turned standard
servers into systems that mimicked I.B.M.'s expensive mainframes, Big Blue fought back. After legal
action failed to fend off the pipsqueak, I.B.M. resorted to a bear hug: it bought Platform in July
for \$150 million. And then it promptly terminated the innovative product.

Despite eliminating the Platform threat, I.B.M. still faces the wrath of many in the computer
industry. The Computer and Communications Industry Association, a trade group backed by the likes of
Google, Oracle and Microsoft, described the Platform deal as ``a clear attempt by I.B.M. to purchase
a company solely to foreclose competition in the mainframe marketplace, protecting its cash cow at
the expense of consumers.''

And T3 Technologies, a small company that resold Platform's products and was devastated by I.B.M.'s
move, has filed an antitrust complaint against I.B.M. with regulators at the European Commission.

Platform was not the only potential competitor that drew I.B.M.'s fire. At the same time that it
sued Platform, I.B.M. declined to renew a patent license with Fundamental Software, which also made
mainframe emulation software. As a result, Fundamental sits in limbo with a once-popular product it
cannot sell, hoping that I.B.M. will change its stance.

Also, in 2007, QSGI, which refurbishes mainframes, complained about enduring ``anticompetitive''
practices that blocked it from aiding customers and said the company might exit the business. It now
refuses to discuss the mainframe business publicly, according to a spokesman.

In a paper commissioned by Microsoft examining the alternative mainframe technologies, Walter
F.~Tichy, a professor of computer science at the University of Karlsruhe in Germany, concluded that,
as a result of I.B.M.'s actions, ``customers have been denied the benefits of technological
innovation and must instead pay above-market prices for I.B.M. mainframe solutions and premium wages
for a dwindling mainframe workforce.''

I.B.M. said in a statement that it was confident no competition laws have been violated.

No stranger to controversies about its clout, I.B.M. agreed in 1956 to an antitrust consent decree
after a battle with the Justice Department. Long since expired -- it had to do with accounting
machines, not mainframes -- the decree nevertheless helped companies like Amdahl, Hitachi and
Fujitsu sell computers that could run I.B.M.'s mainframe software, which they licensed from the
company.

By the late 1990s, all of the other mainframe makers decided to abandon the technology because it
was too expensive to keep up with I.B.M.'s custom chips and software.

More recently, Sun Microsystems, Hewlett-Packard and Microsoft have made mostly unsuccessful
attempts to pull mainframe customers away from I.B.M. by creating products that handle similar tasks
but run on servers.

I.B.M. is now negotiating to buy Sun for about \$7 billion, and if the deal were to occur, I.B.M.
would also gain a monopoly on the key storage systems used for mainframes.

Mainframes crunch the data just about every time someone withdraws money from an automated teller
machine, uses a credit card or buys a product from a large retailer.

I.B.M. contends that the continued popularity of mainframes stems from its efforts to modernize the
systems so that they can run more contemporary business software.

However, A.~M. Sacconaghi, an analyst at Sanford C.~Bernstein \& Company, suggests that I.B.M. has
benefited more from the lack of competition than from updated technology.

The growth rate for the amount of mainframe processing power sold each year has fallen during the
last eight years, undermining I.B.M.'s claims of rising interest in the products. I.B.M.'s release
of new mainframe systems has slowed to every 30 to 36 months from every 18 months or so.

Platform argued that by running emulation technology on standard business servers, it had a cheaper,
faster alternative that could meet the needs of smaller businesses, which had been neglected by
I.B.M. The technology let customers run mainframe and server programs on the same hardware, meaning
they could buy less and do more.

``It really was something that the marketplace wanted,'' said Ron Hilton, a former Amdahl engineer
and later the founder and chief technology officer of Platform. Early customers included the
University of Alabama Medical Center, Cascade Natural Gas and Polk County in Iowa.

H.P. liked Platform's concept, and in 2006, it almost bought the company for close to \$200 million.
Just before the deal was to close, however, it fell apart when H.P.'s lawyers discovered letters
from I.B.M. stating that it would refuse to license its mainframe software to Platform.

I.B.M. sued Platform weeks later, accusing it of infringing on I.B.M.'s patents and undermining the
company's immense investment in mainframe technology. Platform responded with a countersuit,
accusing I.B.M. of seeking to eliminate choice in the mainframe market.

It later complained to the European Commission about I.B.M.'s alleged abuse of its dominant position
in the mainframe market. Platform's lawyers considered Europe a more likely place for a legal
victory, since regulators there tend to have more sympathy for antitrust complaints.

Unable to sell products without an I.B.M. software license, Platform fired most of its staff,
keeping five people to pursue the litigation. In November 2007, Platform got a jolt of cash when
Microsoft joined Platform's existing investors, including Intel Capital and Goldman Sachs, to put
\$37 million more into the company, allowing it to rehire staff and work on a fresh product.

But as the legal proceedings dragged on, Platform's investors grew weary. ``We were six to nine
months from getting a new product to market,'' Gregory Handschuh, the former general counsel at
Platform, recalled. ``The investors just didn't have the stomach for fighting a very difficult
case.''

I.B.M., meanwhile, did not want to risk a lengthy European antitrust investigation, Mr.~Handschuh
said.

I.B.M. maintains that it bought Platform for the company's technology and talent, not to kill a
competitive product.

Platform's engineers had explored ways to speed the flow of data between other computers and a
mainframe, adding horsepower to jobs like encryption or data analysis. They also had a view of how
they could extend these designs to include other accelerators, like I.B.M.'s Cell chip.

``It was at a point where we could gain time to market with the technology, and that means money,''
said Mark Anzani, the chief technology officer for I.B.M.'s mainframe business.

While Platform has disappeared, its fight against I.B.M. lives on in a modified form. T3, the
biggest packager of Platform's technology, is carrying on the battle with financial support from
Microsoft.

In January, T3 presented the European Commission with yet another complaint asserting that I.B.M.
had abused a monopoly position in the mainframe market.

``You can just be pushed around by a schoolyard bully, or you can fight them and stand up for what
you believe is right,'' said T3's president, Steven Friedman. ``We still want to provide an option
for the mainframe marketplace.''

\section{A Quandary in Sweden: Criminals in Med School}

A year ago, Sweden's most prestigious medical school found itself in an international uproar after
it unknowingly admitted a student who was a Nazi sympathizer and a convicted murderer, then
scrambled to find a way to expel him.

It is hard to imagine how the case could get any more bizarre. But it has.

The 33-year-old student, Karl Helge Hampus Svensson, having been banished from the medical school of
the Karolinska Institute in Stockholm on the ground that he falsified his high school records, has
now been admitted to a second well-known medical school -- Uppsala, Sweden's oldest university. New
twists in his and another case highlight the difficulties that three of the country's six medical
schools have had in admitting and dismissing students with serious criminal offenses in just the
past two years. The cases resonate far beyond Sweden, raising fundamental questions about who is fit
to become a doctor.

The circumstances of Mr.~Svensson's admission to Uppsala's first-year class -- reported in January
by Swedish news organizations -- are unknown, because none of the officials involved will publicly
discuss his case. He apparently uses an assumed name -- a customary practice for Swedes seeking to
remain anonymous because of a personal threat. Last week, Uppsala officials, responding to concerns
about Mr.~Svensson's admission, said he had not participated in class work, but did not say why.

In another embarrassing twist, a Swedish newspaper reported last month that much of the verdict and
court files regarding Bjorn Soderberg, Mr.~Svensson's murder victim, had been cut out or replaced
with blank pages. The police said they had been unable to find a culprit.

And in still another case, a 24-year-old medical student at Lund University was convicted last April
of raping a 14-year-old boy while he slept. A district court sentenced the student to two years in
prison, but a higher court reduced the sentence to two years' probation and medical therapy.

When the dean at Lund sought to expel the student, a national board that reviews expulsions blocked
the action, saying that although the man had committed a serious crime, he was not considered a
threat to people or property. The decision was then reversed by an administrative court, which
upheld the expulsion; the student did not appeal.

In contrast with the United States, Swedish laws and customs are sympathetic to released offenders,
saying that once they have served their time they should be treated like ordinary citizens. But the
cases raise questions about protecting the rights of patients and fellow medical students and health
care workers.

Entry to Swedish medical schools is highly competitive. At Uppsala, for example, a spokeswoman said
there were 2,603 applicants for the spring semester and just 100 admissions. They include
Mr.~Svensson, who is taking up a taxpayer-financed slot that could have gone to another student.

Indeed, the Uppsala County Council, which runs all government health facilities in the area, says it
will not allow Mr.~Svensson to do any clinical work, which is a critical and mandatory part of
medical school training. That raises questions about how he will complete a degree even if he does
attend classes.

Mr.~Svensson, who has not responded to numerous attempts to reach him over the last year, was
convicted in the 1999 hate murder of a trade union worker and was paroled after serving 6 ½ years of
an 11-year sentence -- a typical penalty for murder in Sweden. He entered Karolinska in fall 2007
while still on probation; he had earned credits for medical school while in prison.

The disclosures about his past proved deeply embarrassing to the institute. Among other things, two
senior faculty members on the admissions committee that interviewed him failed to ask for an
explanation of the six-and-a-half-year gap in his r\'esum\'e, the period he was in prison.

Swedish universities are legally prohibited from conducting background checks on applicants. To
complicate matters, Mr.~Svensson legally changed his surname from Hellekant after his conviction.

In the United States, the chances of a convicted criminal's being admitted to medical school were
reduced in 2002, when the Association of American Medical Colleges' standard application form began
requiring answers to questions about felony convictions. In 2008, questions were also added about
military discharge history and misdemeanor convictions.

Since Mr.~Svensson had done nothing wrong in medical school, Karolinska officials said they were
powerless to expel him until they discovered his falsified records, blaming the Swedish agency
responsible for checking the validity of educational records.

The Swedish medical licensing agency said that it would not allow Mr.~Svensson to practice even if
he earned his medical degree. But because the agency's jurisdiction excludes universities, questions
arose about whether and how medical school officials should inform patients examined by Mr.~Svensson
about his criminal past, and what the patients' responses would be.

Now Uppsala has declined to say whether its admission committee interviewed him or knew of his
criminal past. Nor is there a public explanation of how he earned any high school credits he may
have been lacking in the relatively short period since his expulsion from the Karolinska.

Last year, Karolinska officials, acting in part on a committee's recommendation, exhorted the
Swedish parliament to enact legislation to allow medical schools a freer hand in admitting and
dismissing students in cases where they would have patient contact.

But little has happened.

Sweden's medical student association is divided about whether a convicted murderer should ever be
allowed to become a doctor.

The president of the student group, Yosef Tyson, said his colleagues have been annoyed about the
very limited journalistic coverage of the cases and the complexities of the issues in training
doctors.

``The problem is not that he's been accepted again; the problem is that nothing was done about this
a year ago,'' Mr.~Tyson said in an interview. ``Everyone is in limbo.''

Mr.~Svensson's case has raised the question of whether Swedish universities should now be allowed to
require criminal records as part of a background check of applicants. While such requirements are
rare in Sweden, they do exist for prospective teachers of young children and for applicants to
private companies where personal security is important, like taxi services.

Another concern is the threat he might pose to patients who are immigrants, or their families --
long a target of neo-Nazi vilification. Even as a student, he will have access to electronic medical
records, which could potentially be misused.

In a group interview, five upper-level Uppsala medical students criticized the university for not
holding an open dialogue about Mr.~Svensson's case, and they worried that the affair was harming
their school's reputation -- and, by extension, theirs.

``Not talking to us, telling us, creates a lot of talk and rumors, when we should be focusing on
school,'' said one of the students, Jessica Svefors, 25.

Pontus Andren, 23, said the issue was one of trust. ``If a rapist or a murderer with neo-Nazi
motives can study to become a doctor, that causes a crisis that affects the entire medical
profession,'' he said. ``When you arrive at a hospital or an emergency room, you might not be alert
or even conscious as a patient, and that puts you in a really vulnerable position.''

But another Uppsala student, Karl-Wilhelm Olsson, 23, said that ``the important factor is whether a
person is a risk to another human being, and it's hard to draw a line.''

He added that while there is no law requiring a university to bar prospective students because of a
criminal past, ``a student should be expelled if he or she is viewed as unfit.''

But Gustav Stalhammar, 25, said Mr.~Svensson should be allowed to become a doctor. ``Who is to say
that he might not become a great doctor, even if it in some ways would feel wrong or awkward to have
a murderer for a colleague?'' he asked. ``It is not fair to have preconceptions about his
character.''

Still, the five students said they were upset because they had not had a chance to meet Mr.~Svensson
and judge his motivations. And they are concerned that he has not publicly expressed regret for his
crime.

``A convicted murderer is a risky candidate for the job, especially if you haven't been able to ask
him why he wants to become a doctor,'' Mr.~Andren said.

Ms.~Svefors said the issue went well beyond the Svensson case. ``I am upset that nobody's done
anything with this issue, that no one has shown civil courage to expel him, that those who have
power to do something are just trying to hand off the problem to someone else,'' she said. ``This is
hurting our reputation as future doctors.''

\section{Rescue Plan, With Fine Print, Dazzles Wall Street}

WASHINGTON -- This time President Obama directed some of the stagecraft. This time Treasury
Secretary Timothy F.~Geithner fleshed out the substance of their long-anticipated program to remove
banks' toxic assets and revive the financial system. And this time the reaction was widely positive,
giving the embattled Mr.~Geithner a critically needed boost.

Mr.~Geithner had a lot on the line in Monday's announcement, despite Mr.~Obama's repeated insistence
that his job was safe. The Treasury secretary's highly publicized unveiling six weeks ago of his
plan's ``framework'' had been panned for lacking details, hobbling Mr.~Geithner as he was trying to
build credibility with the financial markets and clout in Washington.

He suffered further political harm last week for not stepping in sooner to block the
taxpayer-supported American International Group's payment of employee bonuses, and the fallout from
that episode -- including a House vote to impose punitive taxes on the recipients of the bonuses --
threatened to scare off the same investors he needed for his plan to work.

Much still needs to be done to execute the plan. It drew criticism in some quarters for not being
aggressive enough in addressing the bad assets weighing down the financial system, and from
Republicans on Capitol Hill for leaving taxpayers too much at risk.

But the stock market's first-day verdict went in the right direction for Mr.~Geithner this time, up
nearly 7 percent and 500 points, in contrast to the precipitous slide after Mr.~Geithner's first
effort, when his inability to explain in any detail how the program would work left Wall Street
jittery about whether the administration had a workable plan.

A Treasury spokeswoman insisted the only difference was that Mr.~Geithner had the time to complete
details so complicated that they amount to creating a new financial system with global reach. But
beyond the substance, the administration also had a more careful plan in place to introduce the
proposal, because neither Mr.~Geithner nor Mr.~Obama could afford another negative review.

``Did we do things differently? It's self-evident that we did,'' Rahm Emanuel, the White House chief
of staff, said in an interview.

With selective leaks to the media for the last several days, the administration had time to explain
the complexities in advance, preparing the financial markets over the weekend for what was coming.
Mr.~Geithner and other administration officials spent days briefing crucial people on Wall Street
and working to line up endorsements from prominent equity fund managers and other private-sector
``validators,'' in particular two leading global investment management firms, BlackRock and Pimco.

In a White House meeting late last week, Mr.~Obama personally admonished administration officials to
join Mr.~Geithner in the plan's public marketing. As for the Treasury secretary, instead of speaking
awkwardly from teleprompters to an audience of financial V.I.P.'s and 17 television cameras, as he
did on Feb.~10, Mr.~Geithner banned cameras on Monday when he met with reporters to release the plan
and answered questions without notes and at length.

The administration also paid close attention to the political climate. With the private sector
increasingly wary of Congressional intervention in the business of those who participate in
government bailout programs, Mr.~Obama substantially dialed back the near endorsement he had given
late last week to the House vote for a confiscatory 90 percent tax on bonuses like those A.I.G.
doled out.

``As a general proposition, you don't want to be passing laws that are just targeting a handful of
individuals,'' he said in an interview on Sunday night on CBS's ``60 Minutes,'' in what
administration officials said was a signal to investors that he understood their concerns about
doing business with the government.

The Treasury Department also increased the amount of equity the government would provide for the
public-private financing package, to as much as \$100 billion, up from some earlier estimates of as
low as \$50 billion. ``The private sector stepped up on the day of the announcement saying they
wanted to participate, which provided a Good Housekeeping seal,'' Mr.~Emanuel said. ``And the second
thing was that the presentation was organized in a way that the expectations of the program were met
by the actual program.''

For both, he gave credit to Mr.~Geithner and his still-small Treasury team. There was even good news
by Monday evening on the matter of Treasury's much criticized understaffing. The White House
announced the nominations of two Clinton administration veterans to top posts: Neal S.~Wolin to be
Mr.~Geithner's deputy Treasury secretary, and Lael Brainard to be under secretary for international
affairs.

Yet even as Mr.~Geithner tries to focus further on the financial rescue plan, the A.I.G. controversy
will dog him for a second week. On Tuesday he will be on Capitol Hill to answer questions from the
House Financial Services Committee about his handling of the A.I.G. bailout and bonuses.

Chairman Ben S.~Bernanke of the Federal Reserve and William C.~Dudley, Mr.~Geithner's successor as
president of the Federal Reserve Bank of New York, also will testify. Both have shared jurisdiction
over A.I.G. with Treasury since last September, when the government bailed out the global insurance
giant to prevent it from collapsing and taking down the major financial institutions that were among
its customers.

For Mr.~Geithner, the A.I.G. controversy was the latest unexpected setback, after the bad reviews of
the earlier rescue plan announcement and the disclosure during his Senate confirmation process of
some unpaid payroll taxes earlier in the decade.

Nonetheless a new CBS New poll found that a majority of the public -- 54 percent -- say they have
confidence in his ability to handle the nation's financial crisis.

That support masks a partisan divide, as more Democrats express confidence than Republicans do in
Mr.~Geithner. As it happens, Mr.~Geithner used to be a Republican and changed his voter registration
-- to Independent -- only in the late 1990s, when he became under secretary for international
affairs in the Clinton administration.

The poll found that three-fourths of those surveyed said the government should try to recover the
bonus money paid by A.I.G. But the support for recovering the bonus money drops when the method of
recovery is taxation.

\section{Strip-Search of Girl Tests Limit of School Policy}

SAFFORD, Ariz. -- Savana Redding still remembers the clothes she had on -- black stretch pants with
butterfly patches and a pink T-shirt -- the day school officials here forced her to strip six years
ago. She was 13 and in eighth grade.

An assistant principal, enforcing the school's antidrug policies, suspected her of having brought
prescription-strength ibuprofen pills to school. One of the pills is as strong as two Advils.

The search by two female school employees was methodical and humiliating, Ms.~Redding said. After
she had stripped to her underwear, ``they asked me to pull out my bra and move it from side to
side,'' she said. ``They made me open my legs and pull out my underwear.''

Ms.~Redding, an honors student, had no pills. But she had a furious mother and a lawyer, and now her
case has reached the Supreme Court, which will hear arguments on April 21.

The case will require the justices to consider the thorny question of just how much leeway school
officials should have in policing zero-tolerance policies for drugs and violence, and the court is
likely to provide important guidance to schools around the nation.

In Ms.~Redding's case, the United States Court of Appeals for the Ninth Circuit, in San Francisco,
ruled that school officials had violated the Fourth Amendment's ban on unreasonable searches.
Writing for the majority, Judge Kim McLane Wardlaw said, ``It does not require a constitutional
scholar to conclude that a nude search of a 13-year-old child is an invasion of constitutional
rights.''

``More than that,'' Judge Wardlaw added, ``it is a violation of any known principle of human
dignity.''

Judge Michael Daly Hawkins, dissenting, said the case was in some ways ``a close call,'' given the
``humiliation and degradation'' involved. But, Judge Hawkins concluded, ``I do not think it was
unreasonable for school officials, acting in good faith, to conduct the search in an effort to
obviate a potential threat to the health and safety of their students.''

Richard Arum, who teaches sociology and education at New York University, said he would have handled
the incident differently. But Professor Arum said the Supreme Court should proceed cautiously.

``Do we really want to encourage cases,'' Professor Arum asked, ``where students and parents are
seeking monetary damages against educators in such school-specific matters where reasonable people
can disagree about what is appropriate under the circumstances?''

The Supreme Court's last major decision on school searches based on individual suspicion -- as
opposed to systematic drug testing programs -- was in 1985, when it allowed school officials to
search a student's purse without a warrant or probable cause as long their suspicions were
reasonable. It did not address intimate searches.

In a friend-of-the-court brief in Ms.~Redding's case, the federal government said the search of her
was unreasonable because officials had no reason to believe she was ``carrying the pills inside her
undergarments, attached to her nude body, or anywhere else that a strip search would reveal.''

The government added, though, that the scope of the 1985 case was not well established at the time
of the 2003 search, so the assistant principal should not be subject to a lawsuit.

Sitting in her aunt's house in this bedraggled mining town a two-hour drive northeast of Tucson,
Ms.~Redding, now 19, described the middle-school cliques and jealousies that she said had led to the
search. ``There are preppy kids, gothic kids, nerdy types,'' she said. ``I was in between nerdy and
preppy.''

One of her friends since early childhood had moved in another direction. ``She started acting weird
and wearing black,'' Ms.~Redding said. ``She started being embarrassed by me because I was nerdy.''

When the friend was found with ibuprofen pills, she blamed Ms.~Redding, according to court papers.

Kerry Wilson, the assistant principal, ordered the two school employees to search both students. The
searches turned up no more pills.

Mr.~Wilson declined a request for an interview and referred a reporter to the superintendent of
schools, Mark R.~Tregaskes. Mr.~Tregaskes did not respond to a message left with his assistant.

Lawyers for the school district said in a brief that it was ``on the front lines of a decades-long
struggle against drug abuse among students.'' Abuse of prescription and over-the-counter medications
is on the rise among 12- and 13-year-olds, the brief said, citing data from the Office of National
Drug Control Policy.

Given that, the school district said, the search was ``not excessively intrusive in light of
Redding's age and sex and the nature of her suspected infraction.''

Adam B.~Wolf, a lawyer with the American Civil Liberties Union, which represents Ms.~Redding, said
her experience was ``the worst nightmare for any parent.''

``When you send your child off to school every day, you expect them to be in math class or in the
choir,'' Mr.~Wolf said. ``You never imagine their being forced to strip naked and expose their
genitalia and breasts to their school officials.''

In a sworn statement submitted in the case, Safford Unified School District v. Redding, No. 08-479,
Mr.~Wilson said he had good reason to suspect Ms.~Redding. She and other students had been unusually
rowdy at a school dance a couple of months before, and members of the school staff thought they had
smelled alcohol. A student also accused Ms.~Redding of having served alcohol at a party before the
dance, Mr.~Wilson said.

Ms.~Redding said she had served only soda at the party, adding that her accuser was not there. At
the dance, she said, school administrators had confused adolescent rambunctiousness with
inebriation. ``We're kids,'' she said. ``We're goofy.''

The search was conducted by Peggy Schwallier, the school nurse, and Helen Romero, a secretary.
Ms.~Redding ``never appeared apprehensive or embarrassed,'' Ms.~Schwallier said in a sworn
statement. Ms.~Redding said she had kept her head down so the women could not see that she was about
to cry.

Ms.~Redding said she was never asked if she had pills with her before she was searched. Mr.~Wolf,
her lawyer, said that was unsurprising.

``They strip-search first and ask questions later,'' Mr.~Wolf said of school officials here.

Ms.~Redding did not return to school for months after the search, studying at home. ``I never wanted
to see the secretary or the nurse ever again,'' she said.

In the end, she transferred to another school. The experience left her wary, nervous and
distrustful, she said, and she developed stomach ulcers. She is now studying psychology at Eastern
Arizona College and hopes to become a counselor.

Ms.~Redding said school officials should have taken her background into account before searching
her.

``They didn't even look at my records,'' she said. ``They didn't even know I was a good kid.''

The school district does not contest that Ms.~Redding had no disciplinary record, but says that is
irrelevant.

``Her assertion should not be misread to infer that she never broke school rules,'' the district
said of Ms.~Redding in a brief, ``only that she was never caught.''

Ms.~Redding grew emotional as she reflected on what she would have done if she had been told as an
adult to strip-search a student. Dabbing her eyes with a tissue, she said she would have refused.

``Why would I want to do that to a little girl and ruin her life like that?'' Ms.~Redding asked.

\section{In a Volatile Time, Obama Strikes a New Tone for Crisis}

WASHINGTON -- For just under an hour on Tuesday night, Americans saw not the fiery and inspirational
speaker who riveted the nation in his address to Congress last month, or the conversational
president who warmly engaged Americans in talks across the country, or even the jaunty and jokey
president who turned up on Jay Leno.

Instead, in his second prime-time news conference from the White House, it was Barack Obama the
lecturer, a familiar character from early in the campaign. Placid and unsmiling, he was the
professor in chief, offering familiar arguments in long paragraphs -- often introduced with the
phrase, ``as I said before'' -- sounding like the teacher speaking in the stillness of a classroom
where students are restlessly waiting for the ring of the bell.

The session in the East Room came at a volatile moment for the new president as he sought to quell
Democratic misgivings about his ambitious economic agenda and deflect strong Republican opposition.
Speaking past the reporters in the room to the tens of millions of viewers tuning in at home, he
tried to reassure the nation that he could solve the crisis that has gripped the economy for more
than a year.

``We're beginning to see signs of progress,'' he said, calling for a ``renewed confidence that a
better day will come.''

As balky senators from his own party began carving some of the signature proposals out of his
budget, Mr.~Obama signaled that he could compromise in the short term on a middle-class tax cut and
a cap on carbon emissions. But he indicated that he would stand firm on four top priorities,
insisting that Congress make progress in those areas.

``We never expected when we printed out our budget that they would simply Xerox it and vote on it,''
Mr.~Obama said, expressing flexibility about the details as long as his central goals were met.
``The bottom line is that I want to see health care, energy, education and serious efforts to reduce
our budget deficit.''

At a time of anger and anxiety in the country, Mr.~Obama showed little emotion. He rarely cracked a
joke or raised his voice. Even when he declared himself upset over the \$165 million in bonuses paid
this month by the American International Group despite its taxpayer bailout, his voice sounded calm
and unbothered. ``I'm as angry as anybody about those bonuses,'' he said, adding that executives
needed to learn that ``enriching themselves on the taxpayers' dime is inexcusable.''

To a certain extent, Mr.~Obama's demeanor could have been calculated -- an effort, aides said, to
lower the temperature after a supercharged week and nudge the country toward what Mr.~Obama
considers the more pressing issues of fixing the banking system and reviving the economy. Even after
excoriating the A.I.G. executives, he cautioned that ``the rest of us can't afford to demonize every
investor or entrepreneur who seeks to make a profit.''

The only time he seemed irritated came when he was asked why the attorney general of New York,
Andrew M.~Cuomo, seemed to have more success getting A.I.G. executives to return some bonuses than
his own administration. Pressed on why he did not express outrage immediately upon learning of the
bonuses, Mr.~Obama said sharply, ``Well, it took us a couple of days because I like to know what I'm
talking about before I speak.''

Even on one of the most polarizing subjects in American life, race relations, Mr.~Obama deviated
little from the median. Asked about his impact as the first African-American president, he said the
nation experienced ``justifiable pride'' at his inauguration.

``But that lasted about a day,'' he said, in perhaps his only joke of the night. ``Right now the
American people are judging me exactly the way I should be judged and that is, are we taking the
steps to improve liquidity in the financial markets, create jobs, get businesses to reopen, keep
America safe?''

He showed his usual comfort with a wide array of subjects, even as he excluded the nation's big
newspapers from the questioning in favor of a more eclectic mix. He signaled that the new
conservative government in Israel could make achieving a peace deal more difficult. He expressed
patience about dealing with Iran. And he defended his proposal to increase the tax burden on the
wealthy.

This was Mr.~Obama as more enervating than energizing, a reminder of the way he could be in his
early days as a presidential candidate, before he became defined by rapturous crowds.

``He doesn't seem to emote any real urgency or anger,'' said Matthew Dowd, a former Republican
strategist who has often been complimentary of the new president. ``So at times it comes across as a
bit distant and intellectual.''

Joe Trippi, a Democratic consultant, said: ``He said all the right things. But sometimes his
confidence makes him seem flat.''

Still, the news conference came as Mr.~Obama is moving from the flush of his inauguration into the
grind of Congressional negotiations and shifting public opinions.

Over the past month, he has gone from offering near apocalyptic views of the economy to expressing
optimism in its underlying strength.

Appearing on ``60 Minutes,'' he laughed in talking about the problems he faced, leading his
interviewer to ask if he was punch-drunk.

That was not a question that seemed pertinent Tuesday. He did demonstrate an ability to take a
punch, though, and to deliver one. Even beforehand, Republicans criticized his economic agenda,
noting that it would amass more debt than the 43 previous presidents combined. Representative John
A.~Boehner of Ohio, the House Republican leader, called for ``a do-over'' on the budget.

``I just think that this might be the most irresponsible piece of legislation I've seen in my
legislative career,'' Mr.~Boehner said.

Mr.~Obama fired back hours later: ``Some of the Republican critics have a short memory, because as I
recall, I'm inheriting a \$1.3 trillion deficit, annual deficit, from them.''

But his Democratic allies are more immediately troublesome as they busily dismantle his budget and
rewrite it to their own taste. Mr.~Obama, who is scheduled to visit Capitol Hill on Wednesday to
meet with Senate Democrats, indicated that plans for a middle-class tax cut and a market-based cap
on carbon emissions did not necessarily have to be included in the final version of the budget.

Instead, he said he could pursue them independently. The \$787 billion stimulus plan already passed
by Congress authorized a \$400 tax credit for two years, so he has time to find ways to finance it
permanently. And he suggested that Congress could take steps toward alternative energy in the budget
without necessarily incorporating the so-called cap-and-trade system because committees could work
on that separately.

Throughout his time in public life, Mr.~Obama has confronted questions about whether he was too
detached, too analytical, too intellectual. In the campaign, he was as likely to be compared to
Adlai E.~Stevenson as he was to John F.~Kennedy. And if there is a pattern to Mr.~Obama, it is to
lumber through periods like this and then become intense and animated at the first sign of trouble.

Over the long term, Mr.~Obama's calm has served him well, in particular at the critical moment in
the campaign when the economy began its steep slide. ``That is one of the things people like about
him,'' Mr.~Trippi said.

\section{As Chinese Investments in Africa Drop, Hope Sinks}

CONAKRY, Guinea -- Chinese and Guinean workers toil shoulder to shoulder on a sun-blasted
construction site at this crumbling city's edge, building the latest symbol of an old and sturdy
alliance: a \$50 million, 50,000-seat stadium.

This city is littered with such tokens of a friendship that first flowered when Guinea was an
isolated and struggling socialist state in the late 1950s.

But so far Guinea has not gotten what it really wants from the world's fastest growing economy: a
multibillion-dollar deal to build desperately needed infrastructure in exchange for access to the
impoverished nation's vast reserves of bauxite and iron ore.

As global commodity prices have plummeted and several of China's African partners have stumbled
deeper into chaos, China has backed away from some of its riskiest and most aggressive plans,
looking for the same guarantees that Western companies have long sought for their investments:
economic and political stability.

``The political situation is not very stable,'' Huo Zhengde, the Chinese ambassador here, said in an
interview, explaining the country's hesitation to invest billions in Guinea, where a junta seized
power after the death of the longtime president in December. ``The international markets are not
favorable.''

Just a year ago China appeared to be upending the decades-old order in Africa, stepping into the
void left by large Western companies too timid to invest in the continent's resource-rich but
fragile states as the market for copper, tin, oil and timber soared to new heights. In the new
scramble for Africa's riches, China sought a hefty share.

With a no-strings-attached approach and a strong appetite for risk, China seemed to offer Africa a
complete economic and political alternative to the heavily conditioned aid and economic
restructuring that Western countries and international aid agencies pressed on Africa for years,
often with uninspiring consequences. Rising China, seeking friends and resources, seemed to be
issuing blank checks.

Today, China's quest for commodities has not stalled. State-owned companies are bargain-hunting for
copper and iron ore in more stable places like Zambia and Liberia. But Chinese companies are now
driving harder bargains and avoiding some of the most chaotic corners of the continent. African
governments facing falling revenues are realizing that they may still need the West's help after
all.

``We have seen in the recent past Chinese companies wade into countries nobody else would,'' said
Philippe de Pontet, an analyst at the Eurasia Group, a private research firm. ``That may be
changing.''

In 2007 China announced a \$9 billion deal with Congo for access to its giant trove of copper,
cobalt, tin and gold in exchange for developing roads, schools, dams and railways needed to rebuild
a country roughly the size of Western Europe and shattered by more than a decade of war.

But that deal is now in doubt as falling prices have left Congo in a much weaker negotiating
position. It also suddenly finds itself needing the help of the International Monetary Fund, which
has objected to writing off the country's old debt even as Congo takes on what amounts to new
mineral-backed loans from China. Congo's political and ethnic turmoil remains deep, and its economy
is near collapse.

A year ago those factors seemed irrelevant. Chinese companies did not flinch from making deals to
search for oil in the pirate-infested waters off Somalia, or to mine industrial metals in places
like Zimbabwe.

Unlike many Western companies, Chinese state oil companies had no qualms about doing business with
the government of Sudan, which has become an international pariah because of the conflict in Darfur.

China espoused a new model for African investment: mutually beneficial trade between sovereign
nations with none of the meddling so common among Western donors and investors, with their demands
for labor and environmental standards, as well as respect for democracy and human rights.

These policies proved popular among African governments, and trade between Africa and China grew to
more than \$100 billion by 2008, from less than \$10 million in the 1980s. African leaders spoke
openly about China's offer of an alternative to the edicts of Western-dominated institutions like
the International Monetary Fund and the World Bank.

But here in Guinea, which has some of the world's largest deposits of bauxite, an ore needed for
making aluminum, that hope has all but collapsed.

``The Chinese have changed their strategy,'' said Ibrahima Sory Diallo, a senior economist in
Guinea's Ministry of Finance and an advocate for Chinese investment. ``They are not going to inject
\$5 billion into an unstable country in an uncertain market climate.''

French colonists once called Guinea a geological scandal, so rich are its deposits of valuable
minerals. Despite years of mining and billions in profits, Guinea remains one of the poorest and
least developed countries in Africa.

So it is no surprise that Guinea's government, first under Lansana Cont\'e, the strongman who ruled
for 24 years until his death last year, and the junta that replaced him, wanted to tap China's cash
and building expertise.

China's approach to securing minerals in Africa has been to sign agreements to build huge projects
in exchange for minerals. In Angola, this kind of arrangement has guaranteed Chinese access to oil
in Africa's fourth largest oil producer, which is now booming after emerging tattered and broke from
a vicious civil war that lasted decades. Chinese and Angolan officials trumpeted this partnership as
a model for Chinese investment in the continent, a win-win relationship benefiting both countries.

But that formulation has proved problematic in an economic downturn. African governments are now
realizing that these deals are in essence loans against future revenue, and falling prices could
leave them saddled with giant piles of debt.

That is what appears to have happened in Congo. At current prices Congo would struggle to meet the
stringent production targets in the Chinese deal, said Patricia Feeney, executive director of Rights
and Accountability in Development, a Britain-based advocacy group.

``The Congolese have raised expectations so much that they could rely on Chinese and turn their
backs on Western donors, and in the process they have probably managed to alienate people who were
willing to help,'' Ms.~Feeney said.

In Guinea, China has backed away from what Guinean officials portrayed as a done deal to build a
much-needed \$1 billion hydroelectric dam.

``The dam is not a gift; it is an investment,'' said Mr.~Huo, the Chinese ambassador. ``That is what
win-win means.''

Guineans are increasingly suspicious of Chinese investment. Many people see Chinese companies as
being just as exploitative as Western ones, if not more so. After the military took power in
December, it raided Chinese companies suspected of selling fake medicines, but the raids degenerated
into open looting of Chinese businesses, tapping a vein of resentment long suppressed.

Hamidou Cond\'e works bare-chested under the relentless sun, digging a hole for the foundation of a
new hospital being built by a Chinese company, yet another symbol of Chinese-Guinean friendship.

Mr.~Cond\'e, 35, who has two wives and four children, said that he had been digging in the hard rock
with a shovel, pick and ax for two months, but that he had yet to receive any pay from his Chinese
taskmasters.

``We work like slaves,'' Mr.~Cond\'e said. ``And like slaves we are not paid. The Chinese bring
nothing good to Guinea.''

\section{Obama to Set Benchmarks in Fight Against Militants}

WASHINGTON -- President Obama plans to further bolster American forces in Afghanistan and for the
first time set benchmarks for progress in fighting Al Qaeda and the Taliban there and in Pakistan,
officials said Thursday.

In imposing conditions on the Afghans and Pakistanis, Mr.~Obama is replicating a strategy used in
Iraq two years ago both to justify a deeper American commitment and prod governments in the region
to take more responsibility for quelling the insurgency and building lasting political institutions.
``The era of the blank check is over,'' Mr.~Obama told Congressional leaders at the White House,
according to an account of the meeting provided on the condition of anonymity because it was a
private session.

The new strategy, which Mr.~Obama will formally announce Friday, will send 4,000 more troops to
train Afghan security forces on top of the 17,000 extra combat troops that he already ordered to
Afghanistan shortly after taking office, administration and Congressional officials said. But for
now, Mr.~Obama has decided not to send additional combat forces, they said, although military
commanders at one point had requested a total of 30,000 more American troops.

Although the administration is still developing the specific benchmarks for Afghanistan and
Pakistan, officials said they would be the most explicit demands ever presented to the governments
in Kabul and Islamabad. In effect, Mr.~Obama would be insisting that two fractured countries plagued
by ancient tribal rivalries and modern geopolitical hostility find ways to work together and
transform their societies.

American officials have repeatedly said that Afghanistan has to make more progress in fighting
corruption, curbing the drug trade and sharing power with the regions, while they have insisted that
Pakistan do more to cut ties between parts of its government and the Taliban. Mr.~Obama telephoned
President Hamid Karzai of Afghanistan and President Asif Ali Zardari of Pakistan on Thursday to
share the main elements of the strategic review.

Setting benchmarks for Pakistan could be particularly difficult. For years, the United States has
simply paid bills submitted by the Pakistani government for counterterrorism operations, even during
truces when its military was not involved in counterterrorism. Pakistan has resisted linking its aid
to specific performance criteria and officials acknowledged that developing those criteria could be
problematic.

The key elements of Mr.~Obama's plan, with its more robust combat force, its emphasis on training,
and its far-reaching goals, foreshadow an ambitious but risky and costly attempt to unify and
stabilize Afghanistan and Pakistan. Mr.~Obama is unveiling his approach at a time when the conflict
is worsening, the lives of the people are not visibly improving, and the intervention by
American-led foreign powers is increasingly resented.

The goals that Mr.~Obama has settled on may be elusive and, according to some critics, even na\"ive.
Among other things, officials said he planned to recast the Afghan war as a regional issue involving
not only Pakistan but also India, Russia, China, the United Arab Emirates, Saudi Arabia and the
Central Asian states.

His plan envisions persuading Pakistan to stop focusing military resources on its longstanding
enemy, India, so it can concentrate more on battling insurgents in its lawless tribal regions. That
goal may be especially hard to achieve given more than a half century of enmity -- including a
nuclear arms race -- between Pakistan and India.

All told, the 21,000 additional American troops that Mr.~Obama will have authorized almost precisely
matches the original number of additional troops that President George W.~Bush sent to Iraq two
years ago, bringing the overall American deployment in Afghanistan to about 60,000. But Mr.~Obama
avoids calling it a ``surge'' and resisted sending the full reinforcements initially sought by
commanders.

Instead, Mr.~Obama chose to re-evaluate troop levels at a series of specific moments over the next
year, officials said. Approaching the issue in increments may be easier to explain to members of
Mr.~Obama's own party who fear he is getting the country as entangled in Afghanistan as Mr.~Bush did
in Iraq.

The officials said Mr.~Obama planned to frame the American commitment as a counterterrorism mission
aimed at denying havens for Al Qaeda, with three main goals -- training Afghan security forces,
supporting the weak central government in Kabul and securing the population. While the new strategy
will call for expanding Afghan security forces more rapidly, it will not explicitly endorse the
request from American commanders to increase the national police and army to 400,000.

At the same time, Mr.~Obama warned Congressional leaders that he would need more than the \$50
billion in his budget plan for military operations and development efforts. Asked by lawmakers about
the prospect of reconciliation with moderate members of the Taliban, officials said Mr.~Obama
replied that he wanted to sift out hard-core radicals from those who were fighting simply to earn
money.

Senator Carl Levin, Democrat of Michigan and chairman of the Armed Services Committee, emerged from
a briefing with Defense Secretary Robert M.~Gates to declare that in his judgment the
administration's review ``was right on track.'' He said the new strategy would send a significant
number of additional trainers to work with the Afghan National Army and police, part of an overall
strategy to ``transfer responsibilities to the Afghans, both militarily and in terms of economic
development.''

Mr.~Levin, who was part of a bipartisan group that pressed Mr.~Bush to set benchmarks for Iraq two
years ago, embraced the idea of doing the same again for Afghanistan. ``There is a determination to
set some benchmarks for Afghanistan, and that will be incredibly important,'' Mr.~Levin said. ``We
haven't had them in Afghanistan.''

Republicans emerging from briefings at the White House and on Capitol Hill withheld comment. Antonia
Ferrier, a spokeswoman for Representative John A.~Boehner of Ohio, the House Republican leader, said
in a statement that he ``had a constructive meeting at the White House'' and that he would ``reserve
public comment until the president makes his formal announcement.''

Dennis C.~Blair, the administration's director of national intelligence, said the United States
still lacked intelligence about the power structures inside the country and other basic information
necessary for a counterinsurgency campaign. ``We know a heck of a lot more about Iraq on a granular
level than we know about Afghanistan,'' he said.

Speaking with reporters, Mr.~Blair estimated that up to three quarters of the Taliban's rank and
file in Afghanistan could be peeled away from the Taliban's leadership, most of whom are hiding in sanctuaries across the border in Pakistan.

\section{Could It Really Be Him? Yeah, Probably}

IT has been only two months since the Obamas moved into the White House, but here in the nation's
capital, some people are already asking: Have you bumped into your president and first lady yet?

This is no idle question. During the Bush years, Washington got used to a homebody president who
preferred bringing friends into the Executive Mansion to venturing outside it. But these days,
President Obama and his wife, Michelle, are popping up all over this city.

Like basketball? There was Mr.~Obama sitting courtside recently alongside astonished fans at the
Verizon Center as he cheered on the Chicago Bulls in a losing battle against the Washington Wizards.

Enjoy the performing arts? The Obamas have been to the Kennedy Center twice, once to see the Alvin
Ailey dance troupe -- with daughters Malia, 10, and Sasha, 7 -- and once for a musical tribute to
Senator Edward M.~Kennedy.

How about a tasty meal? The Obamas have enjoyed white-tablecloth dining at Equinox, Bobby Van's
Steakhouse, B.~Smith's and Georgia Brown's, and street-corner casual at Ben's Chili Bowl and Five
Guys Burgers and Fries.

They have gone to parent-teacher conferences, school sporting events and visited working-class and
gentrifying communities that have rarely served as stomping grounds for American presidents and
first ladies -- speaking to students at a charter school in a predominantly Hispanic neighborhood,
and worshiping in a black church, among other activities. (The president and friends also tossed a
basketball around at a city-run recreation center.)

``Everywhere you go, you're wondering whether or not you'll run into them,'' said Washington's
mayor, Adrian M.~Fenty, who has lunched with the president and first lady.

Political observers are still debating whether this out-and-about style simply reflects the personal
inclinations of the Obamas or some political calculus (or both). But one thing is clear: No other
modern president has reached out so widely to so many corners of the city, says Doris Kearns
Goodwin, a presidential historian.

That is no surprise to friends of the first family. The Obamas, after all, are city people, former
community organizers who have long felt at home in the urban landscape. Mr.~Obama is the first
president since Richard M.~Nixon to be elected while living in a city neighborhood, in his case,
Chicago's racially and economically diverse Hyde Park. And the Obamas are now eager to explore the
city beyond the White House walls.

``They want their lives not to be confined solely to the White House but rather to become a part of
the urban, vibrant fabric of D.C.,'' Valerie Jarrett, a senior adviser to the president and a close
family friend, said in an interview.

Of course, the social schedule of the president and first lady is also a powerful political tool, a
way to nurture political alliances and to cultivate political narratives. The Obamas can enjoy their
time out on the town while, at the same time, reaping potential dividends by reinforcing their
promise to bring change to Washington and honing an image of openness and accessibility, some
Washington watchers say.

``Let's face it: It's very good for getting re-elected,'' Letitia Baldrige, the White House social
secretary to Jacqueline Kennedy, said of the Obamas' socializing. ``It's a great bank of good will
in which they're making deposits every day.''

Political analysts say that the images of Mr.~Obama hooting and hollering during a basketball game,
eating a hot dog at Ben's Chili Bowl and watching the ballet with his wife and daughters -- pastimes
routinely broadcast to a national audience -- may humanize a politician who is sometimes viewed as
too cerebral and distant.

Dee Dee Myers, a former press secretary for President Clinton, said the outings allow Mr.~Obama to
project ``an accessible glamour'' and to convey a message of hope during bleak economic times. (She
said that even the gregarious Clintons never got out this much.)

``It's very humanizing and very encouraging to people,'' Ms.~Myers said. ``And it's valuable for him
politically.''

Some warn, however, that such a schedule can also carry political risks, particularly if it
undermines the mystique of the presidency, the image of power and command that a president needs to
enact an ambitious agenda. Americans love the idea of the common man in a position of political
power. (Think Jimmy Stewart in ``Mr.~Smith Goes to Washington.'') But they can also lose some
respect if a politician seems too familiar. (Think Jimmy Carter in his cardigan.)

``Every once in a while it's great, but there's a chance of overexposing yourself socially,'' said
Bradley A.~Blakeman, a former aide to President George W.~Bush. ``People scratch their heads and
say, ``Doesn't the president have other things to do, especially in a crisis?' ''

It is certainly a shift from historical precedent. In the 19th century, Washington was mostly viewed
as a humid, uninviting town that presidents escaped from when they could.

In modern times, said Michael Beschloss, a presidential historian, the notion of presidential
engagement with Washington has typically meant ``going to parties in Georgetown or making friends on
Capitol Hill, in other words, engaging with the permanent political establishment here.''

``This is really different,'' said Mr.~Beschloss of the Obamas forays into casual restaurants and
working-class neighborhoods.

The Obamas know that it's different. As the first African-American couple in the White House, they
want to reach beyond the prosperous, predominantly white corridors of Washington.

``We were taught you have to get to know the community you're in, and you have to be a part of that
community,'' Mrs.~Obama said during a visit to Mary's Center, a health clinic that serves a
predominantly Hispanic neighborhood. ``D.C. is our community now, and it's our home.''

The president says he hopes to serve as a bridge in a town long divided between the haves and
have-nots. ``I want to see if we can bring those two Washingtons together,'' Mr.~Obama said in an
interview on the ABC program ``This Week with George Stephanopoulos.''

For ordinary people, the unexpected encounters with the new president and first lady are
astonishing.

Joe Clark, a corporate lawyer who sat near the president at the basketball game, described the
experience as ``surreal.''

``I couldn't believe that he was so accessible that I could literally shake his hand and heckle him
about needing to suit up because his team was losing,'' Mr.~Clark said.

That is not to say that the Obamas can live anything close to a normal life here.

``There really is no going out in public and blending in anymore,'' said Ari Fleischer, a former
press secretary to Mr.~Bush, describing the challenges facing any president. ``That really is one of
the burdens of the job. You go into the restaurant and everyone stands up and applauds. You always
have to shake hands.''

``But when you're sitting at the table, either out of fear of the Secret Service, respect for the
office or old-fashioned decency, people usually leave you alone,'' Mr.~Fleischer said. ``You still
can have a nice meal with your friends.''

The Obamas are clearly scoping out varied restaurants and places to visit.

Mr.~Fenty said it was the president who suggested lunching at Ben's Chili Bowl, a well-known
black-owned restaurant. Eleanor Holmes Norton, Washington's delegate to Congress, said Mrs.~Obama
suggested lunch at B.~Smith's, also black-owned, a Southern-style restaurant near the Capitol.

Mrs.~Obama and her staff also visited Miriam's Kitchen, a soup kitchen, where the first lady bumped
into Bill Richardson, a 46-year-old homeless man. Mr.~Richardson was so stunned that he could barely
stammer thank you as Mrs.~Obama scooped a helping of mushroom risotto onto his plate this month.

``I was expecting some lunch, but this is the president's wife; this is her right here,'' said
Mr.~Richardson, who said he planned to get to a phone as soon as he could. ``I'm going to be like,
'Mom, you're never going to guess who I've seen.' ''

\section{When Stars Twitter, a Ghost May Be Lurking}

The rapper 50 Cent is among the legion of stars who have recently embraced Twitter to reach fans who
crave near-continuous access to their lives and thoughts. On March 1, he shared this insight with
the more than 200,000 people who follow him: ``My ambition leads me through a tunnel that never
ends.''

Those were 50 Cent's words, but it was not exactly him tweeting. Rather, it was Chris Romero, known
as Broadway, the director of the rapper's Web empire, who typed in those words after reading them in
an interview.

``He doesn't actually use Twitter,'' Mr.~Romero said of 50 Cent, whose real name is Curtis Jackson
III, ``but the energy of it is all him.''

In its short history, Twitter -- a microblogging tool that uses 140 characters in bursts of text --
has become an important marketing tool for celebrities, politicians and businesses, promising a
level of intimacy never before approached online, as well as giving the public the ability to speak
directly to people and institutions once comfortably on a pedestal.

But someone has to do all that writing, even if each entry is barely a sentence long. In many cases,
celebrities and their handlers have turned to outside writers -- ghost Twitterers, if you will --
who keep fans updated on the latest twists and turns, often in the star's own voice.

Because Twitter is seen as an intimate link between celebrities and their fans, many performers are
not willing to divulge the help they use to put their thoughts into cyberspace.

Britney Spears recently advertised for someone to help, among other things, create content for
Twitter and Facebook. Kanye West recently told New York magazine that he has hired two people to
update his blog. ``It's just like how a designer would work,'' he said.

It is not only celebrities who are forced to look to a team to produce real-time commentary on daily
activities; politicians like Ron Paul have assigned staff members to create Twitter posts and
Facebook personas. Candidate Barack Obama, as well as President Obama, has a social-networking team
to keep his Twitter feed tweeting.

The famous, of course, have turned to ghostwriters for autobiographies and other acts of
self-aggrandizement. But the idea of having someone else write continual updates of one's daily life
seems slightly absurd.

The basketball star Shaquille O'Neal, for example, is a prolific Twitterer on his account -- The
Real Shaq -- where he shares personal news, jokes and occasional trash talking about opponents with
nearly 430,000 followers.

``If I am going to speak, it will come from me,'' he said, adding that the technology allows him to
bypass the media to speak directly to the fans.

As for the temptation to rely on a team to supply his words, he said: ``It's 140 characters. It's so
few characters. If you need a ghostwriter for that, I feel sorry for you.''

Athletes seem to be purists. Lance Armstrong, only hours after breaking his right collar bone,
tweeted about it, using his left hand. Charlie Villanueva, a forward for the Milwaukee Bucks,
tweeted at halftime from the locker room on March 15 about how ``I gotta step up.'' (His coach,
Scott Skiles, was not pleased with his diversion, but the Bucks did win.)

But for candidates like Mr.~Paul, Twitter is an organizing tool rather than a glimpse behind the
curtain. During the presidential campaign, said Jesse Benton, Mr.~Paul's campaign manager, ``we
assigned a staffer to each social network site. Each was used to generate the same message as a way
to amplify the message and drive people back to our site.''

He said that in rare cases, however, supporters would read more meaning in the online relationship
than was intended. ``On a bunch of social-networking sites, we would get some sincere written notes
that would say 'thank you for letting me be your friend,' '' he recalled.

Many online commentators are appalled at the practice of enlisting ghost Twitterers, but Joseph
Nejman, a former consultant to Ms.~Spears who helped conceive her Web strategy, said there was a
more than a whiff of hypocrisy among critics.

``It's O.K. to tweet for a brand,'' he said, remarking how common it is for companies to have
Twitter accounts, ``but not O.K. for a celebrity. But the truth is, they are a brand. What they are
to the public is not always what they are behind the curtain. If the manager knows that better than
the star, then they should do it.''

In the last couple of months, the Britney Spears Twitter stream has become a model of transparency.
Where the feed once seemed that it was all written personally by Ms.~Spears -- even the blatantly
promotional items about a new album -- lately it can read like a group blog, with some posts signed
``Britney,'' some signed by ``Adam Leber, manager'' and others by ``Lauren.'' That would be Lauren
Kozak, social-media director of britneyspears.com. (Ms.~Spears's management team declined to be
interviewed for this article.)

An unabashed user of ghost Twitterers is Guy Kawasaki, a new-media consultant with more than 80,000
followers, who is full of praise for the two employees who enliven his Twitter feed, often posting
updates while he is on stage addressing a conference.

``Basically, for 99.9 percent of people on Twitter, it is about updating friends and colleagues
about how the cat rolled over,'' he said. ``For a tenth of a percent it is a marketing tool.''

Annie Colbert, a 26-year-old freelance writer from Chicago who is one of Mr.~Kawasaki's ghost
Twitterers, said she judged her performance based on how often her postings for Mr.~Kawasaki are
``retweeted,'' that is, resent by other users of Twitter.

Recently, she said, she had a coup when the actor Ashton Kutcher repeated her post about a YouTube
video showing someone getting high from a ``natural hallucinogen.''

``Facebook is like 'Cheers,' where everyone knows your name,'' she said. ``Twitter is the hipster
bar, where you booze and schmooze people.''

She said she had been considering trying to get other ghost Twitter clients. ``I don't think I could
ghost Twitter for 100 people,'' she said. ``More like 10 clients. I think I would have to get to
know them.''

\pagebreak
\printindex
\end{CJK*}
\end{document}
